\chapter{Continuity of solutions}
\index{Monge--Ampere equation!continuity of solutions}
\index{regularity!BFJ solution}
\index{Kolodziej argument!non-Archimedean}
\index{capacity!and continuity of solutions}
\index{dual complex!support of a measure}
\index{algebraizability!in the BFJ theorem}
\label{chap:bfj-continuity}

A finite-energy solution is an equality of measures; continuity is a
separate regularity theorem.  In the BFJ method it comes from a particularly
rigid hypothesis: the prescribed measure is carried by one finite dual
complex.  The restriction of every psh function to that complex is
continuous, and capacity propagates an inequality known on the support to
the whole Berkovich space.

This chapter first proves that mechanism in the original setting, then lists
the later residue-characteristic and mixed-characteristic extensions with
their exact additional hypotheses.  In particular, none of the theorems
below says that every finite-energy measure has a continuous potential.

The reason for insisting on this from the first page is that the word
``continuity'' has, by now, three different referents in non-Archimedean
pluripotential theory, and they are logically independent.  There is
continuity of the \emph{envelope} $P_\omega f$ of a continuous obstacle
$f$, proved in \Cref{chap:bfj-envelopes} and used everywhere as an input;
there is continuity of the \emph{solution} of a prescribed Monge--Amp\`ere
equation, which is what this chapter is about; and there is continuity of
the \emph{solution map} $\mu\mapsto\varphi_\mu$, a stability statement
which belongs to \Cref{chap:bfj-finite-energy-measures}.  The first is a
statement about a projection operator and needs multiplier or test ideals;
the second is a statement about a single potential and needs capacity; the
third is a statement about a family and needs compactness.  A reader who
merges any two of them will attribute to the theory a theorem it does not
have.

\paragraph{Standing hypotheses.}
Unless a theorem says otherwise---and several of them do---we retain the
original BFJ geometric hypotheses fixed at the beginning of
\Cref{chap:bfj-model-forms} and in force throughout
\Cref{chap:bfj-psh,chap:bfj-compactness,chap:bfj-envelopes,chap:bfj-ma,chap:bfj-capacity,chap:bfj-energy,chap:bfj-differentiability,chap:bfj-finite-energy-measures}:
the field $K$ is complete and discretely valued with residue
characteristic zero and $v_K(\pi)=1$; the variety $X$ is smooth,
connected and projective of dimension $n$ over $K$; and $\omega$ is a
closed model $(1,1)$-form which is semipositive and whose de Rham class
$\{\omega\}$ is ample.  We write, as always,
\[
 V:=(\{\omega\}^n)>0,
 \qquad
 \mathcal M_V(X):=\bigl\{\mu\geq0\ \text{Radon}\ :\
   \mu(X^{\An})=V\bigr\}.
\]
Every Monge--Amp\`ere measure below has total mass $V$, never mass one;
see \eqref{eq:appJ-cl-total-mass}.  \Cref{chap:bfj-envelopes,chap:bfj-capacity}
write $\theta$ for the reference form, and every result quoted from those
chapters is applied with $\theta=\omega$.  The symbol
$\MA_\omega(\varphi)$ denotes the bounded mixed product of
\Cref{thm:bounded-ma-extension} for bounded $\varphi$ and the
nonpluripolar product \eqref{eq:nonpluripolar-ma} for
$\varphi\in\cE^1(X,\omega)$; the two agree where both are defined, by
\Cref{prop:ch26-npp-basic}\,\cref{itm:ch26-npp:bounded}.

The two sections \Cref{sec:ch29-original,sec:ch29-later} deliberately
step outside these hypotheses, and each theorem there restates in full the
setting in which it is asserted.  This is not pedantry: the whole point of
the chapter is that the hypotheses of the continuity theorems are not
interchangeable.

\paragraph{What the reader should learn.}
By the end of the chapter the reader should be able to: state the capacity
estimate \eqref{eq:bfj-kolodziej-sublevel} and say which of its hypotheses
is used at which step; run the non-Archimedean Ko\l odziej iteration
\Cref{lem:ch29-iteration} and produce from it an \emph{explicit} bound on
the oscillation of a solution, \eqref{eq:ch29-kolodziej-bound}, given a
capacity-domination hypothesis on the measure; prove the domination lemma
\Cref{lem:bfj-supported-domination} and explain why both potentials must be
translated by the \emph{same} constant in its proof; state what support on a
dual complex does and does not buy, in particular that it forbids any mass
on a pluripolar set (\Cref{prop:ch29-supported-no-polar-mass}) but does
\emph{not} force the solution to factor through a retraction; write down the
full hypothesis signature \eqref{eq:ch29-signature} of the original
Boucksom--Favre--Jonsson theorem, including the function-field and
completion assumption, and say where each entry is consumed; and quote each
of the four later continuity theorems with its own date, source and extra
hypotheses, without merging them.

\paragraph{Prerequisites.}
From \Cref{chap:bfj-capacity} we use the definition
\eqref{eq:capacity-definition} of $\Capa_\omega$ and its representation
\eqref{eq:ch25-dom-sup}; the positivity statement
\Cref{lem:divisorial-positive-capacity}, which is what makes ``open of
capacity zero'' mean ``empty''; the Choquet continuity
\Cref{prop:ch25-choquet-axioms}\,\cref{itm:ch25-chq:incr}; the sublevel
estimate \Cref{prop:ch25-sublevel-capacity}, with its constant
$C_\omega+nV$ from \eqref{eq:ch25-normalized-bound}, and its corollary
\Cref{cor:ch25-polar-negligible}; the bounded comparison and domination
principles \Cref{cor:bounded-comparison-principle} and
\Cref{prop:bounded-capacity-domination}; and, for the worked examples, the
exact curve computations
\Cref{cor:ch25-curve-capacity-formula,prop:ch25-segment-capacities} and
\Cref{ex:ch25-segment-capacities,ex:ch25-pluripolar-computed}.  From
\Cref{chap:bfj-energy} we use the class \Cref{def:e1}, its order structure
\Cref{thm:ch26-e1-structure}, the nonpluripolar product
\Cref{def:nonpluripolar-ma} with \Cref{prop:ch26-npp-basic}, the mixed
product \Cref{thm:e1-mixed-products}, and the finite-energy comparison
principle \Cref{cor:ch26-e1-comparison}.  From
\Cref{chap:bfj-compactness} we use skeletal equicontinuity
\Cref{thm:bfj-skeletal-equicontinuity} with its explicit modulus
\Cref{cor:ch22-explicit-modulus}, compactness
\Cref{thm:bfj-psh-compactness}, and decreasing regularization
\Cref{thm:bfj-decreasing-regularization}.  From
\Cref{chap:bfj-differentiability} we use orthogonality
\Cref{def:orthogonality-property} and its provenance ledger
\Cref{thm:orthogonality-provenance}; from
\Cref{chap:bfj-finite-energy-measures}, the variational solution
\Cref{prop:variational-maximizer-solves} and the supported existence
theorem \Cref{thm:bfj-supported-cy}.  From
\Cref{chap:topological-tools} we need the continuous-limit clause of
Dini's theorem for upper semicontinuous nets,
\Cref{prop:dini-usc}\,(1).

\section{Orientation}
\label{sec:ch29-orientation}

\subsection{Three continuity statements, and the arrows between them}
\label{sec:ch29-three-continuities}

Write $P_\omega$ for the psh envelope of \Cref{def:ch23-envelope} and, for
$\mu\in\mathcal M_V(X)$, write $\varphi_\mu$ for a solution of
$\MA_\omega(\varphi_\mu)=\mu$ when one exists.  The three statements are:

\begin{description}[style=nextline,leftmargin=1.5em,itemsep=.4\baselineskip]
\item[(C1) Continuity of envelopes.]
$P_\omega f\in C^0(X^{\An})$ for every $f\in C^0(X^{\An})$.  This is
\Cref{thm:bfj-envelope-regularity}, proved by multiplier ideals on models,
and it implies decreasing model regularization
\Cref{thm:bfj-decreasing-regularization}
by \Cref{cor:ch23-decreasing-regularization-pointer}.
\item[(C2) Continuity of solutions.]
$\varphi_\mu\in C^0(X^{\An})$.  This is the subject of the present
chapter.
\item[(C3) Continuity of the solution map.]
$\mu_i\to\mu$ implies $\varphi_{\mu_i}\to\varphi_\mu$ in a specified
topology.  On a fixed dual complex this is
\Cref{prop:fixed-skeleton-stability}: weak measure convergence gives
BFJ convergence of normalized solutions, uniform convergence on that
complex, and convergence of energies.  The maximizing-sequence statement
\Cref{thm:ch28-finite-energy-existence} has a different input.
\end{description}

There is no bare implication among the three.  Under the additional
existence and support hypotheses of this chapter, however, (C1) is an
input to the proof of (C2): through
\Cref{thm:bfj-decreasing-regularization}, the
decreasing net $\varphi_j\searrow\varphi$ of model functions is the object
to which Dini's theorem is applied.  Statement (C1) is also used, through
orthogonality, to produce the solution in the first place.  But (C1) does
\emph{not} imply (C2): the envelope theorem holds for every continuous
obstacle under the standing hypotheses, while
\Cref{ex:ch29-unbounded-not-supported} exhibits, in the same setting, a
measure of finite energy whose unique finite-energy solution is unbounded,
hence not continuous.  Conversely (C2) for one measure says nothing about
(C1).  The fixed-support form of (C3) uses a family of continuous solutions
and adds a compactness argument.  The general finite-energy
maximizing-sequence theorem of \Cref{thm:ch28-finite-energy-existence}
does not require its limiting solution to be continuous.

The mechanism of (C2) is worth stating before any of the estimates, because
it is short.  Suppose $\MA_\omega(\varphi)=\mu$ and let
$\varphi_j\searrow\varphi$ be model functions.  Fix $\varepsilon>0$.
\emph{If} we can find one index $j$ with
\begin{equation}
\label{eq:ch29-mechanism}
 \varphi_j-\varepsilon\ \leq\ \varphi
 \qquad\mu\text{-almost everywhere},
\end{equation}
then a capacity argument---the domination lemma
\Cref{lem:bfj-supported-domination}---upgrades
\eqref{eq:ch29-mechanism} to an inequality \emph{everywhere}, and since
$\varphi_j\geq\varphi$ always, this gives
$0\leq\varphi_j-\varphi\leq\varepsilon$ on all of $X^{\An}$.  Uniform
convergence of continuous functions then makes $\varphi$ continuous.  The
whole difficulty is concentrated in producing \eqref{eq:ch29-mechanism}, and
the device used here is Dini's theorem, which requires a
\emph{fixed compact set} on which $\varphi$ is already known to be
continuous and which carries $\mu$.  Hypothesis \textup{(S5)} supplies
exactly such a set.

\subsection{The hypothesis signature of the original theorem}
\label{sec:ch29-signature}

The theorem historically labelled Theorem~A of \cite{BFJ15} carries six
hypotheses.  We display them once, with the labels used throughout the
chapter, because the entire scope discussion below is an audit of this list:
\begin{equation}
\label{eq:ch29-signature}
\begin{aligned}
 \textup{(S1)}\quad&\text{$K$ complete, discretely valued;}\\
 \textup{(S2)}\quad&\text{residue characteristic $0$;}\\
 \textup{(S3)}\quad&\text{$X$ smooth connected projective, $\omega$ a closed}\\
 &\text{semipositive model form, $\{\omega\}$ ample;}\\
 \textup{(S4)}\quad&\text{$X$ of geometric origin from a
   \emph{one}-dimensional family: there are a curve $B$}\\
 &\text{over the residue field, a point $b\in B$, and a variety
   over $k(B)$}\\
 &\text{whose base change to the fraction field $K$ of
   $\widehat{\cO_{B,b}}$ is $X$;}\\
 \textup{(S5)}\quad&\text{$\mu\geq0$ supported on the dual complex of
   one projective SNC model;}\\
 \textup{(S6)}\quad&\text{$\mu(X^{\An})=V$.}
\end{aligned}
\end{equation}
Hypothesis \textup{(S4)} is the one that disappears from informal accounts.
It is a condition on the \emph{origin} of $X$, not on its analytic
geometry; it says that the pair $(K,X)$ is the completion, at one place,
of a global object over a one-variable function field.  In \cite{BFJ15} it
is used exactly once, in Appendix~A, to prove orthogonality by spreading a
model out over the curve $B$ and applying an absolute volume computation
there.  It was removed---by a different proof of orthogonality, not by
weakening the analytic arguments---in \cite{BGGJKM20}; this is
\Cref{thm:bggjkm-remove-algebraizability} below, and its provenance ledger
is \Cref{thm:orthogonality-provenance}.  Hypothesis \textup{(S5)} is retained in that work's continuous-solution
theorem.  General finite-energy existence in the standing BFJ setting is
\Cref{thm:ch28-finite-energy-existence}; continuity requires a separate
argument.  Skeletal support is sufficient rather than necessary, as the
bounded continuous example in \Cref{ex:ch28-unbounded-dual-energy} shows.
\Cref{sec:ch29-original} is devoted to \textup{(S4)} and \textup{(S5)}.

\newpage
\begin{warning}{Two hypotheses that look alike and are not}
\textup{(S4)} and \textup{(S5)} are both ``geometric'' hypotheses, but concern different objects; the dual complex in \textup{(S5)}
can have dimension greater than one.  \textup{(S4)} concerns a curve $B$ over the
\emph{residue} field, downstairs, and constrains $X$ itself;
\textup{(S5)} concerns a dual complex inside $X^{\An}$, upstairs, and
constrains $\mu$.  The cited 2020 theorem removes \textup{(S4)} while retaining
\textup{(S5)}.  Replacing \textup{(S5)} by finite dual energy alone cannot
give continuity: \Cref{ex:ch29-unbounded-not-supported} is a counterexample.
We never write ``the geometric hypothesis'' without
saying which.
\end{warning}

\subsection{What is imported}
\label{sec:ch29-imported}

Six statements are used here without proof, all of them proved earlier in
this book or, where marked, quoted from the literature with an exact
reference.

\begin{enumerate}
\item\label{itm:ch29-import:reg}
 Decreasing model regularization,
 \Cref{thm:bfj-decreasing-regularization}: every $\omega$-psh function is
 the limit of a decreasing net of $\omega$-psh model functions.  This uses
 \textup{(S1)}, \textup{(S2)} and \textup{(S3)}, through multiplier ideals.
\item\label{itm:ch29-import:equi}
 Skeletal equicontinuity, \Cref{thm:bfj-skeletal-equicontinuity} and
 \Cref{cor:ch22-explicit-modulus}: on a fixed dual complex the normalized
 $\omega$-psh functions form a uniformly bounded, equi-Lipschitz family.
\item\label{itm:ch29-import:cap}
 The capacity toolkit of \Cref{chap:bfj-capacity}: the representation
 \eqref{eq:ch25-dom-sup}, positivity
 \Cref{lem:divisorial-positive-capacity}, the sublevel bound
 \Cref{prop:ch25-sublevel-capacity}, and the bounded comparison principle
 \Cref{cor:bounded-comparison-principle}.
\item\label{itm:ch29-import:e1}
 The finite-energy calculus of \Cref{chap:bfj-energy}: convexity of
 $\cE^1$, the mixed product \Cref{thm:e1-mixed-products}, the maximum
 principle \Cref{thm:ch26-e1-maximum-principle} and the comparison
 principle \Cref{cor:ch26-e1-comparison}.
\item\label{itm:ch29-import:var}
 Existence of a finite-energy solution for a measure supported on a dual
 complex, \Cref{thm:bfj-supported-cy}, which uses orthogonality through
 \Cref{prop:variational-maximizer-solves}.
\item\label{itm:ch29-import:uniq}
 Uniqueness modulo constants of the finite-energy solution,
 \Cref{thm:ch29-uniqueness}, quoted from
 \cite[Proposition~8.2]{BFJ15}.  Its proof uses iterated
 Cauchy--Schwarz and the ample Hodge equality, following Yuan--Zhang.  A complete proof in the model
 regime is given in \Cref{prop:ch29-model-uniqueness}.
\end{enumerate}

These are the analytic inputs to the standing regularity argument.
The later theorems identify their additional imported algebraic and
potential-theoretic inputs separately.  In particular, no result of
\Cref{chap:cld-bedford-taylor} and no local Bedford--Taylor theory is used,
and the Chambert--Loir measure of a merely continuous metric appears only
through the comparison \Cref{thm:bfj-cld-ma-common-locus}.

\subsection{Plan}
\label{sec:ch29-plan}

\Cref{sec:ch29-capacity} proves the two capacity estimates on which
everything rests: the finite-energy sublevel estimate
\Cref{lem:bfj-kolodziej-sublevel}, and the constant-obstacle estimate
\Cref{prop:ch29-constant-obstacle} in the form in which an iteration can
consume it.  \Cref{sec:ch29-kolodziej} isolates the iteration itself,
\Cref{lem:ch29-iteration}, as a statement about a real function of one
variable, and deduces the non-Archimedean Ko\l odziej theorem
\Cref{thm:ch29-kolodziej-linfty} with an explicit constant.
\Cref{sec:ch29-supported} studies measures carried by a dual complex: they
charge no pluripolar set, they pair continuously with normalized psh
functions.  For the segment model studied below, and for finite
divisorial supports, the capacity-domination hypothesis has computable
constants.  \Cref{sec:ch29-regularity} applies
the domination lemma to prove the regularity theorem, and records the uniqueness
statement with its exact provenance.  \Cref{sec:ch29-original} is the
hypothesis audit of \eqref{eq:ch29-signature}.  \Cref{sec:ch29-later}
presents the remaining three later continuity theorems separately, each
with its date, source and extra hypotheses.  \Cref{sec:ch29-worked} carries
the curve, finite-divisorial and boundary examples to explicit numbers, and \Cref{sec:ch29-scope,sec:ch29-dictionary} close the
chapter.

\section{The capacity estimates behind regularity}
\label{sec:ch29-capacity}

\subsection{The sublevel estimate in \texorpdfstring{$\cE^1$}{E1}}
\label{sec:ch29-sublevel}

We retain $V=(\{\omega\}^n)$.  The following estimate is the
non-Archimedean analogue of the elementary Ko\l odziej capacity inequality.
Its bounded case is \Cref{prop:bounded-capacity-domination} of
\Cref{chap:bfj-capacity}, proved there with the same interpolation device;
here both potentials are allowed to be unbounded, as a finite-energy
variational solution may be.

\begin{lemma}\label{lem:bfj-kolodziej-sublevel}
Let $\varphi,\psi\in\mathcal E^1(X,\omega)$, with $\psi\leq0$.  For
$0<t<1$,
\begin{equation}
\label{eq:bfj-kolodziej-sublevel}
 \Capa_\omega\{\varphi<\psi\}
 \leq t^{-n}
 \int_{\{\varphi<(1-t)\psi+t\}}
             \MA_\omega(\varphi).
\end{equation}
\end{lemma}

\begin{proof}
All the measures occurring below are positive of total mass $V$, by
\Cref{thm:e1-mixed-products}, so every measure of a Borel set written
below is finite.  Comparisons of the potentials retain extended-real
order on their polar sets.

\emph{Step 1: the family of competitors.}  By
\eqref{eq:ch25-dom-sup} the capacity of a Borel set $E$ may be computed
as
\begin{equation}
\label{eq:ch29-capacity-as-sup}
 \Capa_\omega(E)
 =\sup\bigl\{\MA_\omega(u)(E)\ :\ u\in\PSH(X,\omega),\ 0\leq u\leq1\bigr\},
\end{equation}
the normalization $0\leq u\leq1$ being obtained from the normalization
$-1\leq u\leq0$ of \eqref{eq:capacity-definition} by a translation, which
changes no Monge--Amp\`ere measure
(\Cref{lem:ch25-admissible-calculus}\,\cref{itm:ch25-adm:transl}).  Fix such
a $u$ and set
\[
 \psi_t:=(1-t)\psi+tu .
\]
Since $\psi\in\cE^1$ and $u$ is bounded, hence in $\cE^1$ by
\Cref{thm:ch26-e1-structure}\,\cref{itm:ch26-e1:bounded}, convexity of the
class (\Cref{thm:ch26-e1-structure}\,\cref{itm:ch26-e1:convex}) gives
$\psi_t\in\cE^1(X,\omega)$.

\emph{Step 2: two inclusions.}  First, $\psi\leq0$ gives $(1-t)\psi\geq\psi$ in extended-real
order, and $tu\geq0$, so $\psi_t\geq\psi$.  This comparison also holds
on $\{\psi=-\infty\}$ and never subtracts two infinite values.  Thus
\begin{equation}
\label{eq:ch29-two-inclusions}
 \{\varphi<\psi\}\ \subseteq\ \{\varphi<\psi_t\}
 \ \subseteq\ \{\varphi<(1-t)\psi+t\},
\end{equation}
the second inclusion because $u\leq1$.  This is the only place where the
hypothesis $\psi\leq0$ is used, and it is used to obtain the \emph{first}
inclusion.  Translating $\psi$ alone changes the strict sets being
estimated, so a common translation is required when preserving an
almost-everywhere comparison between the original potentials.

\emph{Step 3: a lower bound for $\MA_\omega(\psi_t)$.}  Writing
$\omega+\ddc\psi_t=(1-t)(\omega+\ddc\psi)+t(\omega+\ddc u)$ and expanding
the $n$-fold product by multilinearity
(\Cref{thm:e1-mixed-products}, extending
\Cref{cor:bounded-ma-multilinearity}), all $2^n$ terms are positive
measures, and discarding all but the pure $u$-term gives
\begin{equation}
\label{eq:ch29-convexity-lower-bound}
 \MA_\omega(\psi_t)\ \geq\ t^n\,\MA_\omega(u),
\end{equation}
which is \eqref{eq:ma-convex-lower-bound} in the finite-energy category.

\emph{Step 4: comparison.}  The finite-energy comparison principle
\eqref{eq:ch26-e1-comparison}, applied to the pair
$(\varphi,\psi_t)\in\cE^1\times\cE^1$, reads
\[
 \int_{\{\varphi<\psi_t\}}\MA_\omega(\psi_t)
 \ \leq\ \int_{\{\varphi<\psi_t\}}\MA_\omega(\varphi).
\]
Combining this with \eqref{eq:ch29-two-inclusions} and
\eqref{eq:ch29-convexity-lower-bound},
\[
\begin{aligned}
 t^n\MA_\omega(u)\bigl(\{\varphi<\psi\}\bigr)
 &\ \leq\ t^n\MA_\omega(u)\bigl(\{\varphi<\psi_t\}\bigr)
  \ \leq\ \int_{\{\varphi<\psi_t\}}\MA_\omega(\psi_t)\\
 &\ \leq\ \int_{\{\varphi<\psi_t\}}\MA_\omega(\varphi)
  \ \leq\ \int_{\{\varphi<(1-t)\psi+t\}}\MA_\omega(\varphi).
\end{aligned}
\]
Divide by $t^n$ and take the supremum over $u$, using
\eqref{eq:ch29-capacity-as-sup}.
\end{proof}

\begin{remark}[Where each hypothesis is consumed]
\label{rem:ch29-sublevel-hypotheses}
Four things were used, and it is worth separating them.  Convexity of
$\cE^1$ was used to know that $\psi_t$ is a legitimate competitor;
$\psi\leq0$ was used once, in \eqref{eq:ch29-two-inclusions};
multilinearity of the finite-energy mixed product was used in
\eqref{eq:ch29-convexity-lower-bound}; and the comparison principle
\eqref{eq:ch26-e1-comparison}, which in \Cref{chap:bfj-energy} is deduced
from the maximum principle \Cref{thm:ch26-e1-maximum-principle} by
truncation, was used in Step~4.  \emph{No} regularization, no
quasicontinuity and no continuity of envelopes enters
\Cref{lem:bfj-kolodziej-sublevel}; in particular the lemma is available in
any setting in which the finite-energy calculus of \Cref{chap:bfj-energy}
has been established, whatever the residue characteristic.  What is
\emph{not} available in general is the \emph{input} to that calculus.
\end{remark}

\begin{remark}[The shape of the right-hand side]
\label{rem:ch29-shape-of-rhs}
The estimate is useless as stated unless one can say something about the set
$\{\varphi<(1-t)\psi+t\}$, which contains
$\{\varphi<\psi\}$; the whole art consists in choosing $t$ so small that
the enlargement is harmless, while paying only the factor $t^{-n}$.  Two
choices are made in this chapter.  In the domination lemma
\Cref{lem:bfj-supported-domination} the obstacle $\psi$ is bounded and
$t$ is chosen so that the enlargement stays inside a set already known to
be null; there the factor $t^{-n}$ is irrelevant because the right-hand
side is $0$.  In the Ko\l odziej iteration
\Cref{thm:ch29-kolodziej-linfty} the obstacle is a constant, $t$ is a free
parameter, and the factor $t^{-n}$ is the entire content: it is what makes
the iteration converge geometrically.
\end{remark}

\subsection{The constant-obstacle estimate}
\label{sec:ch29-constant-obstacle}

For the iteration we need the estimate in a form in which the parameter
$t$ enters as an additive shift of the level rather than as a
multiplicative interpolation.  Specializing
\eqref{eq:bfj-kolodziej-sublevel} to a constant $\psi\equiv-s$ produces
the level $-(1-t)s+t=-s+t(s+1)$, whose distance from $-s$ grows with
$s$; an iteration fed by that form does not close.  The following
variant, which uses the same interpolation but performs it \emph{inside} the
constant, is the one that does.

\begin{proposition}
\label{prop:ch29-constant-obstacle}
Let $\varphi\in\cE^1(X,\omega)$ and put $\mu:=\MA_\omega(\varphi)$.
Then for all $s\geq0$ and all $0<t\leq1$,
\begin{equation}
\label{eq:ch29-constant-obstacle}
 t^n\,\Capa_\omega\bigl(\{\varphi<-s-t\}\bigr)
 \ \leq\ \mu\bigl(\{\varphi<-s\}\bigr).
\end{equation}
\end{proposition}

\begin{proof}
Let $u\in\PSH(X,\omega)$ with $0\leq u\leq1$, and set
\[
 \chi:=tu-s-t .
\]
Then $tu=(1-t)\cdot0+t\cdot u$ is a convex combination of the
$\omega$-psh functions $0$ and $u$, hence $\omega$-psh by
\Cref{thm:ch21-stability}, and adding the constant $-s-t$ preserves this
(\Cref{prop:ch21-constants}\,\cref{itm:ch21-const:transl}).  The function
$\chi$ is bounded, so $\chi\in\cE^1(X,\omega)$, and
\begin{equation}
\label{eq:ch29-chi-sandwich}
 -s-t\ \leq\ \chi\ \leq\ -s .
\end{equation}

The left-hand inequality in \eqref{eq:ch29-chi-sandwich} gives
$\{\varphi<-s-t\}\subseteq\{\varphi<\chi\}$, and the right-hand one gives
$\{\varphi<\chi\}\subseteq\{\varphi<-s\}$.  Moreover
$\omega+\ddc\chi=(1-t)\omega+t(\omega+\ddc u)$, so multilinearity gives,
exactly as in \eqref{eq:ch29-convexity-lower-bound},
\[
 \MA_\omega(\chi)\ \geq\ t^n\MA_\omega(u).
\]
The comparison principle \eqref{eq:ch26-e1-comparison} applied to the pair
$(\varphi,\chi)$ therefore yields
\[
\begin{aligned}
 t^n\MA_\omega(u)\bigl(\{\varphi<-s-t\}\bigr)
 &\leq t^n\MA_\omega(u)\bigl(\{\varphi<\chi\}\bigr)
  \leq\int_{\{\varphi<\chi\}}\MA_\omega(\chi)\\
 &\leq\int_{\{\varphi<\chi\}}\MA_\omega(\varphi)
  \leq\mu\bigl(\{\varphi<-s\}\bigr).
\end{aligned}
\]
Taking the supremum over $u$ and using \eqref{eq:ch29-capacity-as-sup}
gives \eqref{eq:ch29-constant-obstacle}.
\end{proof}

\begin{remark}[Why the two estimates differ]
\label{rem:ch29-two-obstacle-forms}
The competitor in \Cref{lem:bfj-kolodziej-sublevel} is
$(1-t)\psi+tu$, a convex combination of $\psi$ and $u$; the competitor
in \Cref{prop:ch29-constant-obstacle} is $tu$ shifted by a constant, that
is, a convex combination of $0$ and $u$ shifted.  When $\psi$ is a
constant the two competitors are genuinely different: the first is
$-(1-t)s+tu$, whose difference from $-s$ is $ts+tu\in[ts,ts+t]$,
while the second is $tu-s-t$, whose difference from $-s$ lies in $[-t,0]$.  The gain is exactly
that the enlargement of the sublevel set in
\eqref{eq:ch29-constant-obstacle} is by the \emph{additive} amount $t$,
independent of $s$.  Both estimates are elementary consequences of the
same comparison principle; only the second iterates.
\end{remark}

\subsection{The capacity profile}
\label{sec:ch29-profile}

The iteration in \Cref{sec:ch29-kolodziej} is applied to the function
$s\mapsto\Capa_\omega(\{\varphi<-s\})$, and it uses its monotonicity and decay.  We also record its
right-continuity, although the iteration does not require that property.

\begin{lemma}
\label{lem:ch29-capacity-profile}
Let $\varphi\in\PSH(X,\omega)$ with $\sup_{X^{\An}}\varphi=0$ and define
\begin{equation}
\label{eq:ch29-profile}
 F_\varphi(s):=\Capa_\omega\bigl(\{\varphi<-s\}\bigr),
 \qquad s\geq0 .
\end{equation}
Then $F_\varphi$ is nonincreasing, right-continuous, and satisfies
\begin{equation}
\label{eq:ch29-profile-decay}
 F_\varphi(s)\ \leq\ \frac{C_\omega+nV}{s}\qquad(s>0),
 \qquad\text{in particular}\qquad F_\varphi(s)\xrightarrow[s\to\infty]{}0,
\end{equation}
with $C_\omega$ the constant of \eqref{eq:ch25-normalized-bound}.
Moreover, if $F_\varphi(s_1)=0$ for some $s_1$, then
$\varphi\geq-s_1$ on all of $X^{\An}$.
\end{lemma}

\begin{proof}
Monotonicity is clear, the sets $\{\varphi<-s\}$ being nested and
$\Capa_\omega$ monotone
(\Cref{lem:ch25-inner-outer}\,\cref{itm:ch25-io:mono}).  Each such set is
open, $\varphi$ being upper semicontinuous.

\emph{Right-continuity.}  Fix $s\geq0$.  Then
\[
 \{\varphi<-s\}=\bigcup_{k\geq1}\Bigl\{\varphi<-s-\tfrac1k\Bigr\},
\]
an increasing union of open sets, so
\Cref{prop:ch25-choquet-axioms}\,\cref{itm:ch25-chq:incr} gives
$\Capa_\omega(\{\varphi<-s\})=\lim_kF_\varphi(s+1/k)$.  Since
$F_\varphi$ is monotone, the limit along the sequence $s+1/k$ equals
$\lim_{\sigma\downarrow s}F_\varphi(\sigma)$, which is therefore
$F_\varphi(s)$.

\emph{Decay.}  This is \Cref{prop:ch25-sublevel-capacity}, which is exactly
\eqref{eq:ch29-profile-decay}; it uses the normalization
$\sup\varphi=0$ and nothing else.

\emph{The last assertion.}  The set $\{\varphi<-s_1\}$ is open and has
capacity zero, hence is empty by
\Cref{lem:divisorial-positive-capacity}.
\end{proof}

\begin{remark}[The profile and the iteration endpoint]
\label{rem:ch29-right-continuity}
Right-continuity follows from continuity of capacity along increasing
unions of open sets.  The endpoint in \Cref{lem:ch29-iteration}, however,
is handled by monotonicity alone: if $s_j\uparrow s_\infty$ and
$f(s_j)\to0$, then $0\leq f(s_\infty)\leq f(s_j)$ for every $j$.
No left-continuity assertion is needed.
\end{remark}

The estimate alone does not bound the oscillation of an arbitrary solution.
What BFJ need is a domination statement.

\subsection{The domination lemma}
\label{sec:ch29-domination}

\begin{lemma}\label{lem:bfj-supported-domination}
Let $h\in\PSH(X,\omega)\cap C^0(X^{\An})$ and
$\psi\in\mathcal E^1(X,\omega)$, and put
$\nu=\MA_\omega(\psi)$.
If $h\leq\psi$ $\nu$-almost everywhere, then $h\leq\psi$ everywhere.
\end{lemma}

\begin{proof}
Since $X^{\An}$ is compact and $h$ is continuous, the two numbers
\begin{equation}
\label{eq:ch29-h-bounds}
 m:=\min_{X^{\An}}h,\qquad M:=\max_{X^{\An}}h,
 \qquad \operatorname{osc}(h)=M-m<+\infty,
\end{equation}
are finite.  Fix $\varepsilon>0$ and choose
\begin{equation}
\label{eq:ch29-t-choice}
 0<t<1
 \qquad\text{with}\qquad
 t\bigl(\operatorname{osc}(h)+1\bigr)\ \leq\ \frac\varepsilon2 .
\end{equation}

\emph{The pair to which the estimate is applied.}  Put
\[
 \varphi':=\psi-M+\varepsilon,
 \qquad
 \psi':=h-M .
\]
Both are translates of members of $\cE^1(X,\omega)$, hence lie in
$\cE^1(X,\omega)$ by
\Cref{thm:ch26-e1-structure}\,\cref{itm:ch26-e1:constants}, and
$\psi'\leq0$ by the definition of $M$.  The two potentials have been
translated by the \emph{same} constant $-M$; this is essential and is
discussed in \Cref{warn:ch29-same-translation} below.  Note that
\[
 \{\varphi'<\psi'\}=\{\psi+\varepsilon<h\},
\]
a set unchanged by the common translation, and that
$\MA_\omega(\varphi')=\MA_\omega(\psi)=\nu$, Monge--Amp\`ere measures
being insensitive to additive constants.

\emph{Applying the estimate.}  \Cref{lem:bfj-kolodziej-sublevel}, applied to
the pair $(\varphi',\psi')$, gives
\begin{equation}
\label{eq:ch29-domination-step}
 \Capa_\omega\bigl(\{\psi+\varepsilon<h\}\bigr)
 \ \leq\ t^{-n}\,
   \nu\bigl(\{\varphi'<(1-t)\psi'+t\}\bigr).
\end{equation}
We unwind the set on the right.  Adding $M$ to both sides of the defining
inequality,
\[
 \varphi'<(1-t)\psi'+t
 \iff
 \psi+\varepsilon<h-t(h-M)+t
 \iff
 \psi+\varepsilon<h+t\bigl(M-h+1\bigr).
\]
Since $m\leq h\leq M$ we have $0\leq M-h+1\leq\operatorname{osc}(h)+1$,
so by the choice \eqref{eq:ch29-t-choice},
\[
 \{\varphi'<(1-t)\psi'+t\}
 \ \subseteq\ \Bigl\{\psi+\varepsilon<h+\frac\varepsilon2\Bigr\}
 \ =\ \Bigl\{\psi+\frac\varepsilon2<h\Bigr\}
 \ \subseteq\ \{\psi<h\}.
\]
By hypothesis $\nu(\{h>\psi\})=0$, so the right-hand side of
\eqref{eq:ch29-domination-step} vanishes and
\[
 \Capa_\omega\bigl(\{\psi+\varepsilon<h\}\bigr)=0 .
\]

\emph{Conclusion.}  The set $\{\psi+\varepsilon<h\}=\{h-\psi>\varepsilon\}$
is open: $h$ is continuous and $\psi$ is upper semicontinuous, so
$h-\psi$ is lower semicontinuous with values in
$(-\infty,+\infty]$.  A nonempty open subset of $X^{\An}$ has positive
capacity by \Cref{lem:divisorial-positive-capacity}; hence
$\{\psi+\varepsilon<h\}=\varnothing$, that is,
$h\leq\psi+\varepsilon$ everywhere.  Letting
$\varepsilon\downarrow0$ gives $h\leq\psi$.
\end{proof}

\begin{warning}{Both potentials must be translated by the same constant}
\label{warn:ch29-same-translation}
\Cref{lem:bfj-kolodziej-sublevel} requires its second argument to be
$\leq0$, and the temptation is to arrange this by translating $h$ alone,
so as to write $0\geq h\geq-\operatorname{osc}(h)$.  That is not
legitimate.  The hypothesis of \Cref{lem:bfj-supported-domination} is an
inequality \emph{between} $h$ and $\psi$, and replacing $h$ by
$h-M$ while leaving $\psi$ fixed changes the comparison to
$h-M\leq\psi$ $\nu$-almost everywhere.  When $M>0$, even the resulting
global inequality $h-M\leq\psi$ does not give the desired $h\leq\psi$.  The
proof above therefore translates \emph{both} potentials by $-M$, which
leaves the sets $\{h>\psi\}$ and $\{\psi+\varepsilon<h\}$, the measure
$\nu$, and the hypothesis all unchanged, and buys the normalization
$\psi'\leq0$ for free.  The price is that the constant entering
\eqref{eq:ch29-t-choice} is the \emph{oscillation} of $h$ rather than a
one-sided bound; since $h$ is continuous on a compact space this costs
nothing.  The same point is made, for the bounded statement, in the warning
following \Cref{prop:bounded-capacity-domination}.
\end{warning}

\begin{remark}[Comparison with the bounded domination principle]
\label{rem:ch29-bounded-vs-e1-domination}
The bounded principle \Cref{prop:bounded-capacity-domination} allows
both bounded potentials to be discontinuous and ends with divisorial
order.  Its proof translates both potentials by the same constant.
Here the measured potential $\psi$ may be unbounded; the continuous
minorant $h$ still has finite oscillation, which controls the enlargement
in \eqref{eq:ch29-t-choice}.  The displayed proof concludes by openness
of its strict sublevel set.  These are the exact scopes of the two
arguments; normalization of only one potential is invalid in either.
\end{remark}

\begin{corollary}
\label{cor:ch29-domination-quantitative}
Let $h\in\PSH(X,\omega)\cap C^0(X^{\An})$, let
$\psi\in\cE^1(X,\omega)$ and $\nu=\MA_\omega(\psi)$, and suppose that
for some $\delta>0$,
\[
 \nu\bigl(\{h>\psi+\delta\}\bigr)=0 .
\]
Then $h\leq\psi+\delta$ on all of $X^{\An}$.
\end{corollary}

\begin{proof}
Apply \Cref{lem:bfj-supported-domination} to the pair $(h,\psi+\delta)$;
the translate $\psi+\delta$ lies in $\cE^1$ by
\Cref{thm:ch26-e1-structure}\,\cref{itm:ch26-e1:constants} and has the same
Monge--Amp\`ere measure $\nu$.
\end{proof}

\begin{remark}[Where continuity of $h$ is used]
\label{rem:ch29-continuity-of-h}
Continuity on the compact space supplies the finite bounds for $h$ in
\eqref{eq:ch29-h-bounds}, and it makes
$\{h-\psi>\varepsilon\}$ open in the last step.  In the application $h$
is a model function minus a constant, so both properties hold.
This proof does not establish a domination statement for two arbitrary
unbounded finite-energy potentials.  Nor does it prove that continuity
of a bounded minorant is necessary: the bounded result
\Cref{prop:bounded-capacity-domination} uses divisorial order instead
of openness.  The exact continuous contact-set consequence is
\Cref{cor:ch25-domination-uniqueness}.
\end{remark}

The proof itself uses no skeletal support hypothesis.  That hypothesis
enters in \Cref{sec:ch29-regularity}: it makes an approximation uniform
exactly where the measure lives, after which
\Cref{lem:bfj-supported-domination} propagates the inequality.

\section{The non-Archimedean Ko\l odziej iteration}
\label{sec:ch29-kolodziej}

In the complex theory, Ko\l odziej's method converts a capacity-domination
hypothesis on the right-hand side of a Monge--Amp\`ere equation into an
$L^\infty$ bound for its solution \cite{Kol98,Kol05}.  The conversion is
carried out by an iteration on the sublevel sets, and it is entirely
insensitive to the geometry: only the estimate
\eqref{eq:ch29-constant-obstacle} and elementary real analysis are used.
This section proves the non-Archimedean version with explicit constants.

The reader should be warned at once about its role.  The theorem proved here
is not the step by which \Cref{thm:bfj-supported-regularity} produces a
continuous solution; that step is the domination lemma together with Dini's
theorem, and it produces continuity directly.  What the iteration produces
is boundedness, with an explicit oscillation bound, under capacity
domination.  We do not establish a general implication between this
condition and skeletal support in dimension $\geq2$.  The positive
power in the capacity estimate cannot simply be omitted
(\Cref{ex:ch29-borderline}); this does not make our numerical constant
sharp.  The segment calculation and finite divisorial supports give
concrete applications.

\subsection{The iteration lemma}
\label{sec:ch29-iteration-lemma}

\begin{lemma}
\label{lem:ch29-iteration}
Let $f:[0,+\infty)\to[0,+\infty)$ be nonincreasing.
Suppose there are constants $B>0$ and
$\alpha>0$ such that
\begin{equation}
\label{eq:ch29-iteration-hypothesis}
 t\,f(s+t)\ \leq\ B\,f(s)^{1+\alpha}
 \qquad\text{for all }s\geq0\text{ and all }0<t\leq1 .
\end{equation}
Let $s_0\geq0$ satisfy
\begin{equation}
\label{eq:ch29-iteration-entry}
 f(s_0)\ \leq\ \kappa(B,\alpha):=B^{-1/\alpha}\,2^{-1/\alpha^2}.
\end{equation}
Then $f(s)=0$ for every $s\geq s_0+2$.
\end{lemma}

\begin{proof}
Set $q:=2^{-1/\alpha}\in(0,1)$, so that $q^\alpha=\tfrac12$, and put
$\delta:=\kappa(B,\alpha)$, so that
\begin{equation}
\label{eq:ch29-iteration-delta}
 B\,\delta^\alpha=B\cdot B^{-1}2^{-1/\alpha}=q .
\end{equation}
Define
\[
 t_j:=2^{-j}\ \in(0,1],\qquad
 s_{j+1}:=s_j+t_j\qquad(j\geq0),
\]
so that $s_j\uparrow s_\infty:=s_0+\sum_{j\geq0}2^{-j}=s_0+2$.

\emph{Claim: $f(s_j)\leq\delta q^{\,j}$ for every $j\geq0$.}  For
$j=0$ this is \eqref{eq:ch29-iteration-entry}.  Assume it for $j$.  By
\eqref{eq:ch29-iteration-hypothesis} with $s=s_j$ and $t=t_j$,
\[
 f(s_{j+1})\ \leq\ \frac{B}{t_j}\,f(s_j)^{1+\alpha}
 \ \leq\ B\,2^{\,j}\,\bigl(\delta q^{\,j}\bigr)^{1+\alpha}
 \ =\ \bigl(B\delta^{\alpha}\bigr)\,\delta\,
       q^{\,j}\cdot\bigl(2^{\,j}q^{\,j\alpha}\bigr).
\]
Now $q^{\alpha}=\tfrac12$ gives $q^{\,j\alpha}=2^{-j}$, so the last
bracket equals $1$; and $B\delta^\alpha=q$ by
\eqref{eq:ch29-iteration-delta}.  Hence
$f(s_{j+1})\leq\delta q^{\,j+1}$, which is the claim for $j+1$.

\emph{Conclusion.}  For every $s\geq s_\infty=s_0+2$ and every
$j$ we have $s\geq s_j$.  Monotonicity gives
$0\leq f(s)\leq f(s_j)\leq\delta q^j\to0$, hence $f(s)=0$,
including at the endpoint.  No continuity assumption on $f$ is needed.
\end{proof}

\begin{remark}[The three constants]
\label{rem:ch29-iteration-constants}
The proof used the freedom in \eqref{eq:ch29-iteration-hypothesis} only for
the dyadic values $t=2^{-j}$, and the total displacement $\sum_j2^{-j}=2$
is the source of the additive constant $2$ in the conclusion.  Any
summable choice $t_j=\tau\lambda^{\,j}$ with $0<\lambda<1$ and
$\tau\leq1$ works and produces the additive constant
$\tau/(1-\lambda)$ at the cost of a different $\kappa$; the balance
$\lambda=\tfrac12$ is convenient, not optimal.  The exponent
$\alpha$ enters twice, once through $q=2^{-1/\alpha}$ and once through
$\delta^\alpha$, which is why $\kappa(B,\alpha)$ contains
$2^{-1/\alpha^2}$: as $\alpha\downarrow0$ the entry threshold collapses
superexponentially for this particular choice of constants.  No
optimality of that rate is proved.  \Cref{ex:ch29-borderline} shows only
that at $\alpha=0$ the boundedness conclusion is false.
\end{remark}

\begin{remark}[How the entry point is found]
\label{rem:ch29-iteration-hypotheses}
No continuity hypothesis or limit hypothesis at infinity is needed once
an entry point $s_0$ satisfying \eqref{eq:ch29-iteration-entry} is
supplied.  In the application, \eqref{eq:ch29-profile-decay} both
guarantees such an entry point and gives a sufficient value for it.
\end{remark}

\subsection{The capacity-domination hypothesis}
\label{sec:ch29-domination-hypothesis}

\begin{definition}[Capacity domination]
\label{def:ch29-capacity-domination}
Let $\mu$ be a positive Radon measure on $X^{\An}$, and let
$A>0$ and $\alpha>0$.  We say that $\mu$ is \emph{$(A,\alpha)$-dominated
by capacity} if
\begin{equation}
\label{eq:ch29-capacity-domination}
 \mu(E)\ \leq\ A\,\Capa_\omega(E)^{1+\alpha}
 \qquad\text{for every Borel }E\subseteq X^{\An}.
\end{equation}
\end{definition}

Three remarks fix the meaning of this condition.  First, it depends on
the fixed reference form $\omega$ through $\Capa_\omega$.
No invariance under a change of reference form is asserted; see
\Cref{rem:ch25-independence-of-theta}.  Second, it forces $\mu$ to give no
mass to any set of capacity zero, in particular to any pluripolar set, by
\Cref{cor:ch25-polar-negligible}; so it is at least as strong as the
absolute-continuity condition that appears in the finite-energy theory.
Third, it is not implied by finite energy: \Cref{ex:ch29-borderline} exhibits
a measure of finite energy on $\mathbf P^1$ failing
\eqref{eq:ch29-capacity-domination} for \emph{every} $\alpha>0$.

The following elementary criterion is the one that applies in practice.

\begin{lemma}
\label{lem:ch29-capacity-floor}
Let $\mu$ be a positive Radon measure with $\mu(X^{\An})=V$, and suppose
there is $c>0$ with
\begin{equation}
\label{eq:ch29-capacity-floor}
 \Capa_\omega(\{x\})\ \geq\ c
 \qquad\text{for every }x\in\Supp\mu .
\end{equation}
Then $\mu$ is $(A,\alpha)$-dominated by capacity for every
$\alpha>0$, with
\[
 A=A(c,\alpha):=V\,c^{-(1+\alpha)} .
\]
If $c\geq1$ one may take $A=V$, independently of $\alpha$.
\end{lemma}

\begin{proof}
Let $E$ be Borel.  If $\mu(E)=0$ there is nothing to prove.  Otherwise
$E$ meets $\Supp\mu$: indeed, if $E\cap\Supp\mu=\varnothing$ then
$E\subseteq X^{\An}\setminus\Supp\mu$, which is $\mu$-null by
definition of the support.  Pick $x\in E\cap\Supp\mu$.  Monotonicity of
capacity and \eqref{eq:ch29-capacity-floor} give
$\Capa_\omega(E)\geq\Capa_\omega(\{x\})\geq c$, whence
\[
 \mu(E)\ \leq\ V
 \ =\ V\,c^{-(1+\alpha)}\cdot c^{\,1+\alpha}
 \ \leq\ V\,c^{-(1+\alpha)}\,\Capa_\omega(E)^{1+\alpha},
\]
as asserted.
\end{proof}

\begin{remark}[When the floor is available]
\label{rem:ch29-floor-availability}
Condition \eqref{eq:ch29-capacity-floor} holds automatically when
$\Supp\mu$ is a \emph{finite} set of divisorial points, by
\Cref{lem:divisorial-positive-capacity}: take $c$ to be the minimum of the
finitely many positive numbers $\Capa_\omega(\{x\})$.  For the segment with positive endpoint reference masses in
\Cref{prop:ch29-curve-point-capacity}, it holds with an explicit $c$.
That proposition is not a computation for every curve skeleton.  In dimension $\geq2$ we do \emph{not} assert it for an arbitrary
measure supported on a dual complex; see
\Cref{rem:ch29-floor-open}.  This is the precise sense in which the
Ko\l odziej route and the Dini route to regularity are different arguments
with different hypotheses.
\end{remark}

\subsection{The oscillation theorem}
\label{sec:ch29-linfty}

\begin{theorem}
\label{thm:ch29-kolodziej-linfty}
Assume $n\geq1$.  Let $\varphi\in\cE^1(X,\omega)$ with
$\sup_{X^{\An}}\varphi=0$, and suppose that
$\mu:=\MA_\omega(\varphi)$ is $(A,\alpha)$-dominated by
capacity in the sense of \eqref{eq:ch29-capacity-domination}.  Then
$\varphi$ is bounded, and
\begin{equation}
\label{eq:ch29-kolodziej-bound}
 \operatorname{osc}_{X^{\An}}\varphi
 \ =\ -\inf_{X^{\An}}\varphi
 \ \leq\ \bigl(C_\omega+nV\bigr)\,A^{1/\alpha}\,2^{\,n/\alpha^2}\ +\ 2 ,
\end{equation}
with $C_\omega$ the constant of \eqref{eq:ch25-normalized-bound}.
\end{theorem}

\begin{proof}
Let $F_\varphi$ be the capacity profile \eqref{eq:ch29-profile} and put
\[
 f(s):=F_\varphi(s)^{1/n}=\Capa_\omega\bigl(\{\varphi<-s\}\bigr)^{1/n}.
\]
By \Cref{lem:ch29-capacity-profile} the function $f$ is nonincreasing
and tends to $0$.  It is also right-continuous, although that
property is not used in the iteration.

\emph{Step 1: the iteration hypothesis.}  Let $s\geq0$ and
$0<t\leq1$.  \Cref{prop:ch29-constant-obstacle} gives
\[
 t^n\,F_\varphi(s+t)\ \leq\ \mu\bigl(\{\varphi<-s\}\bigr),
\]
and \eqref{eq:ch29-capacity-domination} applied to the Borel set
$\{\varphi<-s\}$ bounds the right-hand side by
$A\,F_\varphi(s)^{1+\alpha}$.  Hence
\[
 t^n f(s+t)^n\ \leq\ A\,f(s)^{n(1+\alpha)} ,
\]
and extracting $n$-th roots,
\begin{equation}
\label{eq:ch29-iteration-verified}
 t\,f(s+t)\ \leq\ B\,f(s)^{1+\alpha},
 \qquad B:=A^{1/n}.
\end{equation}
This is exactly \eqref{eq:ch29-iteration-hypothesis}.

\emph{Step 2: the entry point.}  By \eqref{eq:ch29-profile-decay},
\[
 f(s)\ \leq\ \Bigl(\frac{C_\omega+nV}{s}\Bigr)^{1/n}
 \qquad(s>0),
\]
so \eqref{eq:ch29-iteration-entry} holds as soon as
\[
\begin{aligned}
 \Bigl(\frac{C_\omega+nV}{s}\Bigr)^{1/n}
 \ \leq\ B^{-1/\alpha}2^{-1/\alpha^2}
 &\iff
 \frac{C_\omega+nV}{s}\ \leq\ B^{-n/\alpha}2^{-n/\alpha^2}\\
 &\iff
 s\ \geq\ (C_\omega+nV)\,B^{\,n/\alpha}\,2^{\,n/\alpha^2}.
\end{aligned}
\]
Since $B^{n}=A$, we may therefore take
\begin{equation}
\label{eq:ch29-entry-point}
 s_0:=(C_\omega+nV)\,A^{1/\alpha}\,2^{\,n/\alpha^2}.
\end{equation}

\emph{Step 3: conclusion.}  \Cref{lem:ch29-iteration} gives $f(s)=0$, hence
$F_\varphi(s)=0$, for every $s\geq s_0+2$.  By the last clause of
\Cref{lem:ch29-capacity-profile}, $\varphi\geq-(s_0+2)$ on all of
$X^{\An}$.  As $\sup\varphi=0$, this is
\eqref{eq:ch29-kolodziej-bound}.
\end{proof}

\begin{corollary}
\label{cor:ch29-kolodziej-floor}
Let $\varphi\in\cE^1(X,\omega)$ with $\sup\varphi=0$ and
$\MA_\omega(\varphi)=\mu\in\mathcal M_V(X)$.  Suppose
\eqref{eq:ch29-capacity-floor} holds with a constant $c>0$ and
$n\geq1$.  Then
\begin{equation}
\label{eq:ch29-kolodziej-floor-bound}
 \operatorname{osc}_{X^{\An}}\varphi
 \ \leq\ \frac{C_\omega+nV}{c}\ +\ 2 .
\end{equation}
\end{corollary}

\begin{proof}
For every $\alpha>0$, \Cref{lem:ch29-capacity-floor} supplies
$A_\alpha=Vc^{-(1+\alpha)}$.  Substitution into
\eqref{eq:ch29-kolodziej-bound} gives
\[
 \operatorname{osc}\varphi
 \leq \frac{C_\omega+nV}{c}
          \left(\frac Vc\right)^{1/\alpha}2^{n/\alpha^2}+2 .
\]
Let $\alpha\to+\infty$.  The left side is independent of $\alpha$,
and both remaining factors tend to one, proving the claim.  If $c\geq1$
this also gives the weaker bound $C_\omega+nV+2$ used in the examples.
\end{proof}

\begin{remark}[What the estimate does and does not give]
\label{rem:ch29-linfty-scope}
\Cref{thm:ch29-kolodziej-linfty} bounds $\varphi$; it does not make
$\varphi$ continuous, and no argument in this chapter deduces continuity
from an $L^\infty$ bound.  In the complex theory the same is true: Ko\l
odziej's continuity theorem needs, beyond the $L^\infty$ estimate, a
stability estimate applied to a family of approximating equations, and no
such family is available here without further hypotheses.  What
\Cref{thm:ch29-kolodziej-linfty} does supply, and what is genuinely useful,
is a bound depending only on the capacity-domination data $(A,\alpha)$ and
on the fixed geometric constants $C_\omega$, $n$, $V$; such a bound is
uniform over families of measures with common data, which is exactly what a
compactness argument consumes.
\end{remark}

\begin{remark}[The estimate is far from sharp]
\label{rem:ch29-linfty-slack}
Even in the most favourable situation, on a curve, the bound
\eqref{eq:ch29-kolodziej-floor-bound} overshoots by an order of magnitude.
\Cref{ex:ch29-curve-solution} computes both sides for a measure on a segment
of length $4$: the true oscillation is $1$ and the bound is $8$.  The
slack has three sources, all of them structural: the passage from a measure
bound to a capacity bound in \Cref{lem:ch29-capacity-floor} discards all
information about how $\mu$ is distributed; the constant $C_\omega$ of
\eqref{eq:ch25-normalized-bound} is a Lipschitz-times-diameter bound rather
than an extremal quantity; and the dyadic iteration pays a factor $2$ in
the exponent at every step.  We do not attempt to improve any of the three.
\end{remark}

\section{Measures supported on a dual complex}
\label{sec:ch29-supported}

\subsection{The class of skeletal measures}
\label{sec:ch29-skeletal-class}

Throughout this section $\mathcal X$ is a projective SNC model of $X$,
$\Delta=\Delta_{\mathcal X}\subseteq X^{\An}$ is its dual complex with the
path metric of \Cref{lem:ch22-path-metric}, and
$p_{\mathcal X}:X^{\An}\to\Delta$ is the retraction.  We write
\begin{equation}
\label{eq:ch29-skeletal-measures}
 \mathcal M_V(\Delta)
 :=\bigl\{\mu\in\mathcal M_V(X)\ :\ \mu(X^{\An}\setminus\Delta)=0\bigr\},
\end{equation}
the set of positive Radon measures of total mass $V$ carried by
$\Delta$.  Since $\Delta$ is compact, $\mathcal M_V(\Delta)$ may be
identified with the set of positive Radon measures of mass $V$ on the
compact metric space $\Delta$.

Three properties of this class are used below.  The first is purely
topological, the second is the reason the class is visible to
pluripotential theory at all, and the third is what makes the variational
argument of \Cref{chap:bfj-finite-energy-measures} converge.

\begin{proposition}
\label{prop:ch29-skeletal-compact}
$\mathcal M_V(\Delta)$ is a weakly compact, convex subset of
$\mathcal M_V(X)$.
\end{proposition}

\begin{proof}
Convexity is clear.  The set $\mathcal M_V(X)$ of positive Radon measures
of fixed mass $V$ on the compact Hausdorff space $X^{\An}$ is the image
of the weakly compact set $\mathcal P(X^{\An})$ of
\Cref{prop:measure-compactness} under the homeomorphism
$\nu\mapsto V\nu$, hence is weakly compact; here the identification of
measures with positive functionals is
\Cref{thm:riesz-interface}.  It remains to see that
$\mathcal M_V(\Delta)$ is weakly closed in it.  Let $\mu_i\to\mu$
weakly with $\mu_i\in\mathcal M_V(\Delta)$, and let
$g\in C^0(X^{\An})$ vanish on $\Delta$.  Writing
$g=g^+-g^-$ with $g^\pm$ continuous, nonnegative and vanishing on
$\Delta$, we get $\int g^\pm\rd\mu_i=0$ for every $i$, hence
$\int g^\pm\rd\mu=0$.  Thus $\mu$ annihilates every continuous function
vanishing on the compact set $\Delta$; by Urysohn's lemma and inner
regularity this forces $\mu(X^{\An}\setminus\Delta)=0$.
\end{proof}

\begin{proposition}
\label{prop:ch29-skeleton-not-polar}
Let $v\in\PSH(X,\omega)$ with $v\not\equiv-\infty$.  Then
$v$ is finite at every point of every dual complex; that is,
\[
 \Delta_{\mathcal X}\cap\{v=-\infty\}=\varnothing
\]
for every SNC model $\mathcal X$.
\end{proposition}

\begin{proof}
By \Cref{thm:bfj-skeletal-equicontinuity}, applied to
$v-\sup_{X^{\An}}v$, which is a legitimate normalization because $v$ is
bounded above (\Cref{prop:ch21-sup-on-skeleton}), the restriction
$v|_{\Delta_{\mathcal X}}$ is Lipschitz on each face of
$\Delta_{\mathcal X}$ with a constant depending only on
$\mathcal X$ and $\omega$; in particular it is real valued.  Explicitly,
\eqref{eq:ch22-skeletal-bound} gives
$-C\leq v-\sup v\leq0$ on $\Delta_{\mathcal X}$, so
$v\geq\sup_{X^{\An}}v-C>-\infty$ there.
\end{proof}

\begin{proposition}
\label{prop:ch29-supported-no-polar-mass}
Let $\mu\in\mathcal M_V(\Delta)$ and let
$P\subseteq X^{\An}$ be $\sigma$-pluripolar in the sense of
\Cref{def:ch25-set-classes}, that is, contained in a countable union of polar sets $\{v_j=-\infty\}$ with
$v_j\in\PSH(X,\omega)\setminus\{-\infty\}$.  Then $\mu(P)=0$; indeed
$P\cap\Delta=\varnothing$, so $P$ is contained in a $\mu$-null Borel
set.
\end{proposition}

\begin{proof}
By \Cref{prop:ch29-skeleton-not-polar} each $\{v_j=-\infty\}$ is disjoint
from $\Delta$, hence so is their union and so is $P$.  Therefore
$P\subseteq X^{\An}\setminus\Delta$, a Borel set of $\mu$-measure zero
by \eqref{eq:ch29-skeletal-measures}.
\end{proof}

\begin{remark}[Two different statements about pluripolar sets]
\label{rem:ch29-two-polar-statements}
\Cref{prop:ch29-supported-no-polar-mass} is stronger, and much cheaper,
than the corresponding statement for a general finite-energy measure.  For
an arbitrary $\varphi\in\cE^1$, the measure $\MA_\omega(\varphi)$ gives
no mass to pluripolar sets by
\Cref{prop:ch26-npp-basic}\,\cref{itm:ch26-npp:polar}, and that is a theorem
about the nonpluripolar product.  Here the conclusion is set-theoretic: a
skeletal measure does not merely fail to charge a pluripolar set, it lives
on a compact set disjoint from every one of them.  In particular no
truncation or layer argument is needed to integrate an unbounded
$\omega$-psh function against $\mu$; by
\Cref{thm:bfj-skeletal-equicontinuity} the integrand is bounded on
$\Delta$.  This is the mechanism behind
\Cref{lem:skeletal-measure-finite-energy}, and it is worth isolating,
because boundedness of the restriction supplies the integral directly,
even when the potential is unbounded elsewhere.
\end{remark}

\begin{proposition}
\label{prop:ch29-skeletal-pairing}
Let $\mu\in\mathcal M_V(\Delta)$, and let $C=C(\mathcal X,\omega)$ be
the equicontinuity constant of \Cref{thm:bfj-skeletal-equicontinuity}.  Then
for every $\varphi\in\PSH(X,\omega)$ with $\sup_{X^{\An}}\varphi=0$,
\begin{equation}
\label{eq:ch29-pairing-bound}
 -CV\ \leq\ \int_{X^{\An}}\varphi\,\rd\mu\ \leq\ 0 ,
\end{equation}
and the map
$\varphi\mapsto\int\varphi\,\rd\mu$ is continuous on the compact space of
normalized $\omega$-psh functions for the BFJ topology
\Cref{def:ch22-bfj-topology}.  Moreover, if $\mu_i\to\mu$ weakly in
$\mathcal M_V(\Delta)$, then
$\int\varphi\,\rd\mu_i\to\int\varphi\,\rd\mu$ uniformly in $\varphi$
over that space.
\end{proposition}

\begin{proof}
By \eqref{eq:ch22-skeletal-bound} the restriction $\varphi|_\Delta$
satisfies $-C\leq\varphi\leq0$, and $\mu$ is carried by $\Delta$ of
mass $V$; this is \eqref{eq:ch29-pairing-bound}.

Continuity is immediate from the definition: by
\Cref{def:ch22-bfj-topology} the BFJ topology is the topology of uniform
convergence on every dual complex, so $\varphi_i\to\varphi$ implies
$\varphi_i|_\Delta\to\varphi|_\Delta$ uniformly, and since $\mu$ has
finite mass carried by $\Delta$,
\[
 \Bigl|\int\varphi_i\,\rd\mu-\int\varphi\,\rd\mu\Bigr|
 \leq V\,\lVert\varphi_i-\varphi\rVert_{C^0(\Delta)}
 \longrightarrow0 .
\]
Compactness of the normalized psh space is
\Cref{thm:bfj-psh-compactness}.

For the last assertion, the restrictions of normalized psh functions form a
uniformly bounded, equi-Lipschitz, hence relatively compact subset
$\mathcal K\subseteq C^0(\Delta)$ by
\Cref{cor:ch22-arzela-ascoli}.  Weak convergence of measures is uniform on
a compact set of test functions.  Indeed, given $\varepsilon>0$, cover
$\overline{\mathcal K}$ by finitely many balls of radius
$\varepsilon/(3V)$ centred at $g_1,\dots,g_N\in C^0(\Delta)$; extend
each $g_k$ to $\widetilde g_k\in C^0(X^{\An})$ by the Tietze extension
theorem, which changes no integral since all the measures involved are
carried by $\Delta$; choose $i_0$ with
$\bigl|\int\widetilde g_k\rd\mu_i-\int\widetilde g_k\rd\mu\bigr|
<\varepsilon/3$ for all $k$ and all $i\geq i_0$; and for
$\varphi$ normalized psh, picking $k$ with
$\lVert\varphi|_\Delta-g_k\rVert_\infty\leq\varepsilon/(3V)$, estimate
\[
 \Bigl|\int\varphi\,\rd\mu_i-\int\varphi\,\rd\mu\Bigr|
 \leq V\frac{\varepsilon}{3V}+\frac\varepsilon3+V\frac{\varepsilon}{3V}
 =\varepsilon .\qquad
\]
\end{proof}

\subsection{The capacity floor on a curve}
\label{sec:ch29-curve-floor}

We verify \eqref{eq:ch29-capacity-floor} for the segment with positive
endpoint reference masses below.  Use the notation of
\Cref{sec:ch25-curves}: $C$ is a smooth projective geometrically connected
$K$-curve, $\theta=\omega$ is determined on a split regular semistable
model $\mathcal C$, and all node residue degrees are one.
The graph $\Gamma=\Delta_{\mathcal C}$ carries the modulus lengths,
possibly with intermediate vertices, and
$p:C^{\An}\to\Gamma$ the retraction, and
$\mu_0=\MA_\omega(0)$ is the curvature measure of $\omega$, a positive
measure of mass $V=\deg L$ supported on the vertices.

\begin{proposition}
\label{prop:ch29-curve-point-capacity}
Let $\Gamma$ be a segment of length $\ell>0$ joining the two divisorial
points $x_{E_0},x_{E_1}$, with
$\mu_0=m_0\delta_{x_0}+m_1\delta_{x_\ell}$ and $m_0,m_1>0$.  Write
$x_c$ for the point of $\Gamma$ at distance $c$ from $x_{E_0}$.
Then, for $0<c<\ell$,
\begin{equation}
\label{eq:ch29-curve-point-capacity}
 \Capa_\omega\bigl(\{x_c\}\bigr)
 =\Capa_\omega\bigl(p^{-1}(\{x_c\})\bigr)
 =\min\Bigl\{m_0,\frac1c\Bigr\}+\min\Bigl\{m_1,\frac1{\ell-c}\Bigr\},
\end{equation}
and consequently
\begin{equation}
\label{eq:ch29-curve-floor}
 \inf_{x\in\Gamma}\Capa_\omega(\{x\})
 \ \geq\ \min\Bigl\{m_0,\frac1\ell\Bigr\}
      +\min\Bigl\{m_1,\frac1\ell\Bigr\}\ >\ 0 .
\end{equation}
\end{proposition}

\begin{proof}
The second equality in \eqref{eq:ch29-curve-point-capacity} is the last
display of \Cref{prop:ch25-segment-capacities}, with $a=b=c$.  For the
first, monotonicity gives
$\Capa_\omega(\{x_c\})\leq\Capa_\omega(p^{-1}(\{x_c\}))$, so only the
lower bound is at issue, and it follows by inspecting the maximizer
constructed in the proof of \Cref{prop:ch25-segment-capacities}.  With
$\alpha:=\min\{m_0,1/c\}$ and $\beta:=\min\{m_1,1/(\ell-c)\}$, that
maximizer is
\[
 f(x)=\begin{cases}
   -1+\alpha(c-x),&0\leq x\leq c,\\
   -1+\beta(x-c),&c\leq x\leq\ell,
 \end{cases}
\]
and $u:=f\circ p$ lies in $\mathcal U_\omega$.  By
\eqref{eq:ch25-segment-measure} its Monge--Amp\`ere measure is
\[
 \MA_\omega(u)=(m_0-\alpha)\delta_{x_0}
   +(\alpha+\beta)\delta_{x_c}
   +(m_1-\beta)\delta_{x_\ell},
\]
whose middle term is an atom \emph{at the single point} $x_c$, not merely
on the tube over it.  Hence
$\Capa_\omega(\{x_c\})\geq\MA_\omega(u)(\{x_c\})=\alpha+\beta$, which is
the asserted value.

For \eqref{eq:ch29-curve-floor}, let $0<c<\ell$.  Then $c\leq\ell$ gives
$1/c\geq1/\ell$, whence
$\min\{m_0,1/c\}\geq\min\{m_0,1/\ell\}$; and $\ell-c\leq\ell$ gives
$\min\{m_1,1/(\ell-c)\}\geq\min\{m_1,1/\ell\}$.  Adding the two
inequalities and using \eqref{eq:ch29-curve-point-capacity} gives the bound
for interior points.  At the two endpoints,
\eqref{eq:ch25-cap-endpoint} gives the values
$m_0+\min\{m_1,1/\ell\}$ and $\min\{m_0,1/\ell\}+m_1$, each of which is
at least $\min\{m_0,1/\ell\}+\min\{m_1,1/\ell\}$.  Positivity is clear
since $m_0,m_1,\ell>0$.
\end{proof}

\begin{remark}[The bound \eqref{eq:ch29-curve-floor} is not the exact
infimum]
\label{rem:ch29-floor-not-sharp}
The inequality in \eqref{eq:ch29-curve-floor} is genuinely an inequality.
For $\ell=4$ and $m_0=m_1=1$, the right-hand side is
$\tfrac14+\tfrac14=\tfrac12$, while the exact infimum of
$c\mapsto\min\{1,1/c\}+\min\{1,1/(4-c)\}$ over $[0,4]$ is $1$, attained
at the midpoint: on $(0,1]$ the function equals $1+1/(4-c)\geq\tfrac54$,
on $[1,3]$ it equals $1/c+1/(4-c)$, which is convex with minimum $1$ at
$c=2$, on $[3,4)$ it equals $1/c+1\geq\tfrac54$, and at the two
endpoints it equals $\tfrac54$.  For $\ell=4$ and
$m_0=m_1=\tfrac12$ the infimum is instead $\tfrac34$, attained at the endpoints and not at
the midpoint: the location of the infimum is not a function of the
combinatorics alone, and \eqref{eq:ch29-curve-floor} with
$\lambda=1/\ell$ is sharp in neither case.  For the purposes of
\Cref{lem:ch29-capacity-floor} only a positive lower bound is needed, and
\eqref{eq:ch29-curve-floor} supplies one; in a concrete example the exact
infimum may of course be used, as in
\Cref{ex:ch29-curve-solution}.
\end{remark}

\begin{corollary}
\label{cor:ch29-curve-floor}
Let $\mu\in\mathcal M_V(\Gamma)$ with $\Gamma$ as in
\Cref{prop:ch29-curve-point-capacity}.  Then
\eqref{eq:ch29-capacity-floor} holds with
$c=\min\{m_0,1/\ell\}+\min\{m_1,1/\ell\}$, and
\Cref{thm:ch29-kolodziej-linfty} applies to any solution of
$\MA_\omega(\varphi)=\mu$.
\end{corollary}

\begin{proof}
$\Supp\mu\subseteq\Gamma$, and \eqref{eq:ch29-curve-floor} bounds the
capacity of every point of $\Gamma$ from below by $c$.  Now apply
\Cref{lem:ch29-capacity-floor}.
\end{proof}

\begin{remark}[Metric, not combinatorial]
\label{rem:ch29-floor-metric}
For a fixed pair $m_0,m_1$ and $\ell\geq\max\{1/m_0,1/m_1\}$, the
constant supplied by \Cref{cor:ch29-curve-floor} is $c=2/\ell$, so the
$(A,\alpha)$-domination constant $A=Vc^{-(1+\alpha)}$ of
\Cref{lem:ch29-capacity-floor} grows like $(\ell/2)^{1+\alpha}$, and the
oscillation bound \eqref{eq:ch29-kolodziej-bound} grows like a power of
$\ell$.  Thus the floor \eqref{eq:ch29-curve-floor} degrades as the edge
length grows, that is, as the node of the special fibre becomes thicker.
This is the same phenomenon recorded in
\Cref{ex:ch25-segment-capacities}: the capacity of a divisorial point is not
bounded below uniformly over all models, and any estimate built on
\eqref{eq:ch29-capacity-floor} carries the metric of the skeleton in its
constant.  A general dual complex has finitely many faces of finite
diameter, so a floor, if it exists, is again a metric quantity; the present argument provides no such general higher-dimensional
construction or uniform floor.
\end{remark}

\begin{remark}[The floor in higher dimension is not established here]
\label{rem:ch29-floor-open}
For $x$ a divisorial point, \Cref{lem:divisorial-positive-capacity} gives
$\Capa_\omega(\{x\})>0$; for $x$ a rational quasimonomial point of
$\Delta_{\mathcal X}$ one may pass, by
\Cref{thm:skeleton-subdivision}, to a toroidal blowup on which $x$ becomes
divisorial, and conclude the same.  Neither argument is uniform in $x$:
the first has no explicit constant, and in the second the constant depends
on the subdivision, hence on the denominator of $x$.  On a curve the
uniformity is recovered by the explicit computation
\eqref{eq:ch29-curve-point-capacity}, whose mechanism---restriction to a
one-dimensional convex programme---does not give a higher-dimensional proof.  The Burgos Gil--Sombra
example in \Cref{rem:ch23-curve-specialness} obstructs the analogous
retraction-positivity argument; it does not prove failure of every
possible capacity-floor construction.  We therefore state
\eqref{eq:ch29-capacity-floor} as a \emph{hypothesis} for $n\geq2$, verify
it in the two cases named in \Cref{rem:ch29-floor-availability}, and do not
assert it in general.  Nothing in the proof of
\Cref{thm:bfj-supported-regularity} depends on it.
\end{remark}

\subsection{What support on a dual complex does not mean}
\label{sec:ch29-support-not}

Hypothesis \textup{(S5)} of \eqref{eq:ch29-signature} is used in this
chapter for exactly one purpose: to provide a \emph{fixed compact set},
independent of the approximation index, on which a decreasing net of model
functions converges uniformly and which carries the measure.  It is worth
listing four statements that it does \emph{not} provide, each of which is
occasionally assumed in informal accounts.

\begin{enumerate}
\item\label{itm:ch29-notsupp:factor}
 It does not say that the solution factors through the retraction
 $p_{\mathcal X}$.  See \Cref{warn:ch29-support-not-skeletal} below.
\item\label{itm:ch29-notsupp:model}
 It does not say that the solution is a model function.  Already for a
 curve measure, \Cref{ex:ch29-curve-solution} gives a non-PL solution.
 For rational classes over $\mathbf C(\!(\pi)\!)$,
 \cite[Theorem~A(ii)]{BJGreen24} gives divisorial Green functions that
 are not $\mathbf Q$-PL but remain $\mathbf R$-PL.  Theorem~B of that
 source gives a finite atomic example on $\mathbf P^3$ that is not
 invariant under any SNC retraction and is not a real model function;
 see \Cref{ex:ch28-divisorial-green}.
\item\label{itm:ch29-notsupp:floor}
 It does not, for $n\geq2$, give the capacity floor
 \eqref{eq:ch29-capacity-floor}; see \Cref{rem:ch29-floor-open}.
\item\label{itm:ch29-notsupp:cofinal}
 It is preserved under model domination: if
 $\mathcal X'$ dominates $\mathcal X$, then
 $\Delta_{\mathcal X}\subseteq\Delta_{\mathcal X'}$ as sets, so the
 hypothesis is preserved when passing to a \emph{larger} complex, but a
 measure carried by $\bigcup_{\mathcal X}\Delta_{\mathcal X}$, the union
 over all SNC models, need not be carried by any single one.  The union is dense in $X^{\An}$
 (\Cref{prop:snc-cofinal}), but a dense carrying locus need not carry
 every measure.  The condition is therefore not vacuous.  It is weaker
 than carrying by one complex; the ray example
 \Cref{ex:ch29-unbounded-not-supported} separates them.
\end{enumerate}

Item \cref{itm:ch29-notsupp:cofinal} deserves emphasis because it is the
limitation of applying this particular Dini--domination proof on
successively larger compact sets.
The following observation, which costs nothing, makes it quantitative.

\begin{proposition}
\label{prop:ch29-one-complex}
Let $(\mathcal X_m)_{m\geq1}$ be a sequence of SNC models with
$\Delta_m:=\Delta_{\mathcal X_m}$ increasing, and let
$\mu\in\mathcal M_V(X)$ be carried by
$\bigcup_m\Delta_m$ but by no single $\Delta_m$.  Put
\[
 \eta_m:=\mu\bigl(X^{\An}\setminus\Delta_m\bigr).
\]
Then $\eta_m>0$ for every $m$, while $\eta_m\downarrow0$.
Consequently, for every $\varphi\in\PSH(X,\omega)$, every decreasing net
$(\varphi_j)$ of model functions with $\varphi_j\searrow\varphi$, every
$\varepsilon>0$ and every $m$, Dini's theorem supplies an index
$j$ with $\varphi_j-\varepsilon\leq\varphi$ on $\Delta_m$, and the
argument only bounds the possible failure set by
$X^{\An}\setminus\Delta_m$, of measure $\eta_m>0$.
No $m$ makes this support bound zero; the actual inequality may hold
on a larger set, or everywhere.
\end{proposition}

\begin{proof}
That $\eta_m>0$ for every $m$ is precisely the hypothesis that $\mu$
is carried by no single $\Delta_m$.  The sets $\Delta_m$ increase, so
the numbers $\eta_m$ decrease, and
\[
 \lim_m\eta_m=V-\lim_m\mu(\Delta_m)
  =V-\mu\Bigl(\bigcup_m\Delta_m\Bigr)=V-V=0
\]
by the increasing-union identity for finite measures, the sets
$\Delta_m$ being compact and hence Borel.

For the last clause: $\varphi|_{\Delta_m}$ is continuous by
\Cref{thm:bfj-skeletal-equicontinuity} and the restrictions
$\varphi_j|_{\Delta_m}$ decrease to it, so
\Cref{prop:dini-usc}\,(1) applies on the compact set $\Delta_m$ and gives
such an index $j$.  The inequality so obtained is an inequality on
$\Delta_m$, and $\mu(X^{\An}\setminus\Delta_m)=\eta_m>0$.
\end{proof}

\begin{remark}[Why a diagonal choice does not repair this]
\label{rem:ch29-no-diagonal}
Given $\varepsilon>0$, \Cref{prop:ch29-one-complex} allows one to choose
$m$ with $\eta_m<\varepsilon$ and then $j=j(m)$ with
$\varphi_j-\varepsilon\leq\varphi$ on $\Delta_m$.  The resulting
support argument bounds the failure measure by
$\eta_m<\varepsilon$, but does not establish a null failure set.
The stated domination lemma \Cref{lem:bfj-supported-domination}
requires exactly such a null set: its proof concludes that a certain open set has
capacity zero, and ``capacity at most $\varepsilon$'' does not make an
open set empty.  One could hope to run the argument with
\Cref{cor:ch29-domination-quantitative} instead, but the hypothesis there is
again an exact null set.  This is the content of
\Cref{ex:ch29-dini-localized}.
\end{remark}

\section{From a finite-energy solution to a continuous one}
\label{sec:ch29-regularity}

\subsection{The regularity theorem}
\label{sec:ch29-regularity-theorem}

\begin{theorem}\label{thm:bfj-supported-regularity}
Let $\mu$ have mass $V$ and be supported on a fixed SNC dual complex
$\Delta$.  If $\varphi\in\mathcal E^1(X,\omega)$ satisfies
$\MA_\omega(\varphi)=\mu$, then $\varphi$ is continuous.
\end{theorem}

\begin{proof}
Choose, by \Cref{thm:bfj-decreasing-regularization}, a decreasing net
$(\varphi_j)_{j\in J}$ of $\omega$-psh model functions with
$\varphi_j\searrow\varphi$ pointwise on $X^{\An}$.  No normalization or
global bound on this net is needed: the domination lemma normalizes each
continuous minorant internally by its own maximum.

\emph{Step 1: uniform convergence on $\Delta$.}  By
\Cref{prop:ch29-skeleton-not-polar} the restriction
$\varphi|_\Delta$ is real valued, and by
\Cref{thm:bfj-skeletal-equicontinuity} it is Lipschitz on each face, hence
continuous, on the compact polyhedral complex $\Delta$.  The restrictions
$\varphi_j|_\Delta$ are continuous and decrease pointwise to
$\varphi|_\Delta$.  Dini's theorem in the form
\Cref{prop:dini-usc}\,(1)---or already
\Cref{prop:dini-nets}, the functions here being continuous---applies to the
net $(\varphi_j|_\Delta)$ on the compact space $\Delta$ and gives
uniform convergence.  Thus, for every $\varepsilon>0$ there is
$j_1=j_1(\varepsilon)\in J$ such that
\begin{equation}
\label{eq:dini-on-support}
                    \varphi_j-\varepsilon\leq\varphi
                    \qquad\text{on }\Delta,\qquad j\geq j_1 .
\end{equation}
Directedness of $J$ is what turns the finite subcover in the proof of
\Cref{prop:dini-usc} into a single index; no countability is used, and no
sequence is extracted.

\emph{Step 2: transport to a $\mu$-almost-everywhere statement.}  By
hypothesis $\mu(X^{\An}\setminus\Delta)=0$, so
\eqref{eq:dini-on-support} says precisely that
\[
 h:=\varphi_j-\varepsilon\ \leq\ \varphi
 \qquad\mu\text{-almost everywhere},\qquad j\geq j_1 .
\]
Here $h\in\PSH(X,\omega)\cap C^0(X^{\An})$: it is a model function minus a
constant, hence $\omega$-psh by
\Cref{prop:ch21-constants}\,\cref{itm:ch21-const:transl} and continuous by
\Cref{lem:model-function-affine}.

\emph{Step 3: propagation.}  Apply \Cref{lem:bfj-supported-domination} with
this $h$ and with $\psi=\varphi\in\cE^1(X,\omega)$, noting that
$\nu=\MA_\omega(\varphi)=\mu$ is exactly the measure with respect to which
the almost-everywhere hypothesis has just been verified.  The conclusion is
$h\leq\varphi$ on all of $X^{\An}$, that is
\[
                  0\leq\varphi_j-\varphi\leq\varepsilon
                  \qquad\text{on }X^{\An},\qquad j\geq j_1 ,
\]
the left inequality because the net decreases to $\varphi$.

\emph{Step 4: conclusion.}  Since $\varepsilon>0$ was arbitrary, the net
of continuous functions $(\varphi_j)$ converges to $\varphi$ uniformly
on $X^{\An}$.  A uniform limit of continuous functions on a topological
space is continuous, so $\varphi\in C^0(X^{\An})$.
\end{proof}

\begin{corollary}
\label{cor:ch29-uniform-model-limit}
In the situation of \Cref{thm:bfj-supported-regularity}, the solution
$\varphi$ is a uniform limit of $\omega$-psh model functions.
If $\omega=c_1(L,\lVert\cdot\rVert_0)$ for a semipositive model
metric on a line bundle $L$, the metric
$\rme^{-\varphi}\lVert\cdot\rVert_0$ is continuous semipositive in the
sense of Zhang
\cite{Zhang1995}, and its Chambert--Loir measure in the sense of
\cite{CL06} equals $\mu$.
\end{corollary}

\begin{proof}
The proof of \Cref{thm:bfj-supported-regularity} produced a net
$(\varphi_j)$ of $\omega$-psh model functions converging to
$\varphi$ uniformly.  Inductively choose $j_k$ dominating both $j_{k-1}$
and an index valid for $\varepsilon=1/k$.  Then
$(\varphi_{j_k})_{k\geq1}$ is a \emph{sequence} of $\omega$-psh model
functions converging uniformly to $\varphi$.  Uniform convergence of model
potentials gives uniform convergence of their associated model
metrics when the displayed line-bundle reference metric exists, so its
limit is continuous semipositive.  Its
Chambert--Loir measure is the weak limit of the model measures
$\MA_\omega(\varphi_{j_k})$ by
\Cref{prop:ch24-continuous-uniform}\,\cref{itm:ch24-cu:unif}, and that limit
is $\MA_\omega(\varphi)=\mu$ by \Cref{thm:bounded-ma-extension}.
\end{proof}

\begin{remark}[Why Dini is used only on the support]
\label{rem:ch29-two-dinis}
Only \Cref{prop:dini-usc}\,(1), the version with a continuous limit, is
used, and only on $\Delta$.  Applying it on all of $X^{\An}$ would assume
the conclusion, since $\varphi$ is not yet known to be continuous there.
No preliminary normalization is required: for each fixed $j$ the function
$h=\varphi_j-\varepsilon$ is already continuous and $\omega$-psh, and the
proof of \Cref{lem:bfj-supported-domination} translates both $h$ and
$\varphi$ by the finite number $\max h$.  The argument uses neither a
uniform oscillation bound for the approximants nor metrizability of
$X^{\An}$.
\end{remark}

\begin{remark}[The theorem is about one measure at a time]
\label{rem:ch29-one-measure}
The Dini index $j_1(\varepsilon)$ depends on the potential and its chosen
approximation.  This proof gives uniform approximation of that solution,
rather than a common modulus on all of $X^{\An}$ for a family of solutions.
On any fixed dual complex a common Lipschitz modulus is already provided
by \Cref{cor:ch22-explicit-modulus}.  Fixed-skeleton solution convergence
and energy convergence are \Cref{prop:fixed-skeleton-stability}.
The oscillation estimate \Cref{thm:ch29-kolodziej-linfty} is a different
kind of bound.
\end{remark}

\begin{center}
\fbox{\begin{minipage}{0.91\textwidth}
\textbf{Where the hypotheses enter.}
Orthogonality is used to obtain a variational solution; it is not used in
the preceding regularity proof.  Support on one dual complex is used in
\eqref{eq:dini-on-support}; finite energy supplies comparison and full mass.
Decreasing model approximation is supplied by the BFJ regularization
theorem.  Continuity of envelopes is a separate input; in later extensions
it combines with differentiability of non-Archimedean volume to yield
orthogonality.  These roles should not be merged.
\end{minipage}}
\end{center}

The following table records, step by step, what the proof of
\Cref{thm:bfj-supported-regularity} consumed.

\begin{center}
\small
\begin{tabular}{p{.10\textwidth}|p{.38\textwidth}|p{.42\textwidth}}
\toprule
\textbf{Step}&\textbf{Uses}&\textbf{Does not use}\\
\midrule
1
 &\Cref{thm:bfj-decreasing-regularization}, hence
  \textup{(S1)}--\textup{(S3)} through multiplier ideals;
  \textup{(S5)}, to have a compact $\Delta$;
  \Cref{thm:bfj-skeletal-equicontinuity};
  \Cref{prop:ch29-skeleton-not-polar};
  \Cref{prop:dini-usc}\,(1)
 &\textup{(S4)}; orthogonality; any capacity estimate\\
2
 &\textup{(S5)} again, now as the null-set statement
  $\mu(X^{\An}\setminus\Delta)=0$
 &compactness of the psh space\\
3
 &\Cref{lem:bfj-supported-domination}, hence
  \eqref{eq:bfj-kolodziej-sublevel},
  \eqref{eq:ch26-e1-comparison} and
  \Cref{lem:divisorial-positive-capacity}; \textup{(S6)} through the mass
  identity
 &\textup{(S4)}; orthogonality; the Ko\l odziej iteration\\
4
 &nothing beyond uniform convergence
 &---\\
\bottomrule
\end{tabular}
\end{center}

\subsection{Uniqueness}
\label{sec:ch29-uniqueness}

Existence and regularity have now been separated.  Uniqueness is a third,
independent statement, and it is the only ingredient of
\Cref{cor:bfj-original-continuous-solution} that this book does not prove in
full generality.  We first prove it where it can be proved by the
intersection theory already available, and then quote the general statement
with its exact provenance.

\begin{proposition}
\label{prop:ch29-model-uniqueness}
Let $\mathcal X$ be a regular projective model of $X$ on which
$\omega$ is determined, assume $n\geq1$, and let $f,g$ be model functions determined on
$\mathcal X$ such that both $\omega+\ddc f$ and $\omega+\ddc g$ are
$\mathcal X$-positive, that is, represented on $\mathcal X$ by
relatively \emph{ample} classes.  If
\[
 \MA_\omega(f)=\MA_\omega(g),
\]
then $f-g$ is constant.
\end{proposition}

\begin{proof}
Write $u:=f-g$, a model function determined on $\mathcal X$, and let
$D\in\Div_0(\mathcal X)_{\mathbf R}$ be the vertical divisor with
$\varphi_D=u$, which exists and is unique as a representative of the actual function
$u$.  By \Cref{thm:ch20-ddc-exact}, only the divisor representing
$\ddc u$ is determined modulo $\mathbf R\,\mathcal X_s$.  Put
$\alpha:=\{\omega+\ddc f\}_{\mathcal X}$ and
$\beta:=\{\omega+\ddc g\}_{\mathcal X}$ in $N^1(\mathcal X/R)$; both are
relatively ample by hypothesis, in particular relatively nef.

\emph{Step 1: telescoping.}  As classes,
\[
 (\omega+\ddc f)^n-(\omega+\ddc g)^n
 =\ddc u\wedge\sum_{k=0}^{n-1}
   (\omega+\ddc f)^{\wedge k}\wedge(\omega+\ddc g)^{\wedge(n-1-k)} ,
\]
an identity of signed atomic measures obtained by expanding the difference
of two $n$-th powers of the sum $\omega+\ddc g+\ddc u$; every term is a
mixed model measure in the sense of \Cref{def:mixed-model-measure}.
Integrating $u$ against it and using
$\MA_\omega(f)=\MA_\omega(g)$,
\begin{equation}
\label{eq:ch29-uniqueness-telescope}
 \begin{aligned}
 0&=\int u\,\bigl[\MA_\omega(f)-\MA_\omega(g)\bigr]\\
  &=\sum_{k=0}^{n-1}
    \int u\,\ddc u\wedge(\omega+\ddc f)^{\wedge k}
      \wedge(\omega+\ddc g)^{\wedge(n-1-k)} .
 \end{aligned}
\end{equation}

\emph{Step 2: each summand vanishes.}  By
\eqref{eq:model-dirichlet-negativity}, every summand in
\eqref{eq:ch29-uniqueness-telescope} is $\leq0$, the forms
$\omega+\ddc f$ and $\omega+\ddc g$ being semipositive.  A sum of
nonpositive numbers vanishes only if each vanishes; hence, taking the term
$k=0$,
\[
 \int u\,\ddc u\wedge(\omega+\ddc g)^{\wedge(n-1)}=0 .
\]

\emph{Step 3: the equality case.}  By Step~1 of the proof of
\Cref{prop:model-integration-by-parts}, the left-hand side of the last
display is the vertical intersection number
$\bigl(D^2\cdot\beta^{\,n-1}\bigr)$.  Thus
$\bigl(D^2\cdot\beta\cdots\beta\bigr)=0$ with $\beta$ relatively ample,
and the equality clause of \Cref{thm:ch20-mixed-hodge} gives
$D\in\mathbf R\,\mathcal X_s$.  By the sign normalization
\eqref{eq:ch20-sign} the model function of $\mathcal X_s$ is the constant
$1$, so $u=\varphi_D$ is constant.
\end{proof}

\begin{remark}[Why ampleness cannot be relaxed by approximation]
\label{rem:ch29-ample-not-approximable}
The equality clause of \Cref{thm:ch20-mixed-hodge} is asserted only for
relatively ample classes: for merely nef $\alpha_i$ the inequality
\eqref{eq:ch20-mixed-hodge} is obtained by letting
$\varepsilon\downarrow0$ in
$\alpha_i+\varepsilon\mathcal A$, and a strict inequality
$\bigl(D^2\cdot(\alpha+\varepsilon\mathcal A)^{n-1}\bigr)<0$ for every
$\varepsilon>0$ is compatible with equality in the limit.  Consequently
\Cref{prop:ch29-model-uniqueness} does not extend to semipositive model
metrics by an approximation argument, and \emph{a fortiori} not to
finite-energy potentials.  This is not a defect of the exposition: the
general uniqueness statement is genuinely a different theorem, proved by
different means.
\end{remark}

\begin{theorem}
\label{thm:ch29-uniqueness}
Assume the standing hypotheses.  If
$\varphi,\psi\in\cE^1(X,\omega)$ satisfy
$\MA_\omega(\varphi)=\MA_\omega(\psi)$, then
$\varphi=\psi+c$ for a constant $c\in\mathbf R$, including on their
common polar set.  In particular a solution of $\MA_\omega(\varphi)=\mu$ in
$\cE^1(X,\omega)$ is unique modulo constants, and a continuous
semipositive metric solving $c_1(L,\lVert\cdot\rVert)^n=\mu$ is unique up
to a positive scalar.
\end{theorem}

\begin{proof}[Imported input and its architecture]
This is \cite[Proposition~8.2]{BFJ15}, not merely the continuous
uniqueness clause of its Theorem~A.  Equal measures first give
$I_\omega(\varphi,\psi)=0$ and vanishing of all relevant finite-energy
Dirichlet terms.  The proof in that source iterates Cauchy--Schwarz to
transfer this vanishing to the reference form, including for unbounded
finite-energy potentials.  It then chooses a sufficiently positive ample
model representative on higher models and applies the ample Hodge
equality, following Yuan--Zhang, to obtain equality at all divisorial
points up to a constant.  Divisorial determination gives the global
conclusion.  The transfer and ample replacement are essential:
nondegeneracy is not inferred for the original merely nef mixed classes.
These are quoted inputs, also described in
\Cref{rem:ch28-uniqueness-gap}; the elementary relatively ample
model case is proved above.  Neither orthogonality nor \textup{(S5)}
is used.
\end{proof}

\begin{remark}[Three rigidity statements and their inputs]
\label{rem:ch29-uniqueness-ladder}
\Cref{thm:ch29-uniqueness} treats equality of measures for arbitrary
finite-energy potentials and concludes equality modulo constants.
\Cref{cor:ch26-gauge-nondegenerate}\,\cref{itm:ch26-nondeg:both}
treats zero gauge for continuous potentials and concludes equality of the
actual functions.  \Cref{cor:ch25-domination-uniqueness} treats two
continuous potentials whose measures are each carried by their contact
set and uses bounded domination alone.  These statements have different
normalizations and hypotheses; no strict hierarchy among their
hypotheses is asserted.  On a curve, continuous uniqueness also follows
from the maximum principle for a harmonic difference, as explained in
\Cref{thm:curves-arbitrary-field-continuity}.
\end{remark}

\begin{warning}{Support is not skeletal invariance}
\label{warn:ch29-support-not-skeletal}
The support hypothesis in \Cref{thm:bfj-supported-regularity} supplies the
compact set on which Dini convergence is used.  It does not say that the
solution factors through the retraction to that set.  Boucksom--Jonsson
construct a finite atomic measure supported on a dual complex whose solution
is invariant under no SNC retraction \cite[Theorem~B]{BJGreen24}.  Thus the
continuity proof above is sharp about the role of support and must not be
rewritten as a skeletal-factorization argument.
\end{warning}

\begin{corollary}\label{cor:bfj-original-continuous-solution}
Assume $K$ is complete discretely valued with residue characteristic zero,
$X$ is smooth connected projective, $\omega$ is a closed semipositive
model form with ample class and
has the orthogonality property, and $\mu$ is a positive measure of mass
$V$ supported on one SNC dual complex.  Then there is a continuous
$\omega$-psh solution of $\MA_\omega(\varphi)=\mu$, unique modulo
constants.
\end{corollary}

\begin{proof}
The variational theorem gives an $\mathcal E^1$-solution, and
\Cref{thm:bfj-supported-regularity} makes it continuous.  Uniqueness
is Proposition~8.2 of \cite{BFJ15}, ultimately the non-Archimedean
Calabi--Yau uniqueness theorem of Yuan--Zhang.

In detail: \Cref{lem:skeletal-measure-finite-energy} shows that a measure
supported on $\Delta$ has finite dual energy, and
\Cref{thm:bfj-supported-cy}---whose proof runs the variational argument on
the compact normalized psh space of
\Cref{thm:bfj-psh-compactness}, with the pairing controlled by
\Cref{prop:ch29-skeletal-pairing} and the equation extracted by
\Cref{prop:variational-maximizer-solves}---produces
$\varphi\in\cE^1(X,\omega)$ with $\MA_\omega(\varphi)=\mu$.  That step,
and only that step, uses the orthogonality property.
\Cref{thm:bfj-supported-regularity} then makes $\varphi$ continuous,
using \textup{(S5)} a second time and orthogonality not at all, and
\Cref{cor:ch29-uniform-model-limit} exhibits it as a uniform limit of model
functions.  Uniqueness modulo constants is \Cref{thm:ch29-uniqueness}.
\end{proof}

\section{The original theorem: function field, completion, and support}
\label{sec:ch29-original}

\subsection{The two geometric hypotheses, stated exactly}
\label{sec:ch29-two-geometric}

The historical Theorem~A of \cite{BFJ15} assumed in addition that $X$ was
algebraizable: it descended from a smooth projective variety over a
one-variable function field whose completion is $K$.  That assumption was
used in Appendix~A of BFJ to prove orthogonality.  It was not used in the
capacity argument above.

Because this hypothesis is the one most often dropped in summaries, we
restate it in the form in which it is used, and separate it from the support
hypothesis with which it is sometimes conflated.

\begin{definition}[Algebraizability over a one-variable function field]
\label{def:ch29-algebraizable}
Let $K$ be complete discretely valued with residue field $k$.  A smooth
projective $K$-variety $X$ is \emph{algebraizable over a one-variable
function field}, or \emph{of geometric origin from a one-dimensional family
over $k$}, if there exist
\begin{enumerate}
\item\label{itm:ch29-alg:curve}
 a normal one-dimensional variety $B$ over $k$ and a closed point
 $b\in B$, so that $\cO_{B,b}$ is a discrete valuation ring with
 fraction field the one-variable function field $F:=k(B)$;
\item\label{itm:ch29-alg:completion}
 an isomorphism $\widehat{\cO_{B,b}}\xrightarrow{\ \sim\ }K^\circ$ of
 rings, so that $K$ is the \emph{completion} of $F$ at the place $b$;
\item\label{itm:ch29-alg:descent}
 a projective $F$-variety $Y$ and an isomorphism
 $Y\otimes_FK\simeq X$ over $K$.
\end{enumerate}
If in addition a line bundle $L$, respectively a closed $(1,1)$-form
$\omega$, on $X$ is induced by base change from $Y$, respectively from
a model of $Y$ over $\cO_{B,b}$, one says that $(X,L)$, respectively
$(X,\omega)$, is algebraizable.
\end{definition}

This is \Cref{def:ch23-geometric-origin} with $d=1$, stated again here
because the case $d=1$ is the only one occurring in \cite{BFJ15}, and
because the three clauses play distinguishable roles.

\begin{remark}[What each clause of \Cref{def:ch29-algebraizable} is for]
\label{rem:ch29-algebraizability-clauses}
Clause \cref{itm:ch29-alg:curve} supplies the one-dimensional base over
which a model of $X$ may be spread out.  After compactifying $B$ and the
spread-out data, one obtains the proper global setting in which sections
are finite-dimensional and volumes are computable.  Clause
\cref{itm:ch29-alg:completion} is the one whose omission does the most
damage: it is not enough that $X$ be defined over \emph{some} field with a
discrete valuation, nor that $K$ contain a function field.  What is
required is that $K^\circ$ be the completed local ring of $B$ at $b$,
so that model metrics, vertical divisors and intersection numbers upstairs
correspond exactly to their counterparts over $\cO_{B,b}$, and so that
Berkovich analytification over $K$ agrees with the analytification of the
completion.  Clause \cref{itm:ch29-alg:descent} finally transports the
variety.  Only with all three does the volume of a metrized line bundle on
$X^{\An}$ become an arithmetic-geometric quantity computable on a
projective $k$-variety of dimension $n+1$, which is what
\cite[Appendix~A]{BFJ15} needs.
\end{remark}

\begin{example}[The standard instance and descent to a specified field]
\label{ex:ch29-algebraizability-instance}
Let $k$ be a field of characteristic zero, $B=\mathbf A^1_k$ with
coordinate $t$, and $b=\{t=0\}$.  Then
$F=k(t)$, $\cO_{B,b}=k[t]_{(t)}$ and
$K=k(\!(t)\!)$, which is complete and discretely valued with residue field
$k$.  Every $X$ obtained by base change from a smooth projective
$k(t)$-variety is algebraizable in the sense of
\Cref{def:ch29-algebraizable}.

Descent to this specified field $k(t)$ is an additional requirement,
not a consequence of completeness.  For example, a smooth plane quartic
over $\mathbf C(\!(t)\!)$ with coefficients algebraically independent
over $\mathbf C(t)$ has a generic moduli point: the space of plane
quartics modulo projective coordinates has dimension six, so its moduli
coordinates have positive transcendence degree over $\mathbf C(t)$.
It therefore cannot descend to that specified field.  This example
concerns descent to $\mathbf C(t)$; failure over one specified field is
not by itself a proof of failure of \Cref{def:ch29-algebraizable}, which
allows other curve function fields.  Under the standing hypotheses this costs
nothing today, because \Cref{thm:bggjkm-remove-algebraizability} removes
\textup{(S4)}; in equal characteristic $p$, where the analogous
hypothesis \Cref{def:ch23-geometric-origin} is still needed, it is a real
constraint.
\end{example}

\begin{warning}{Algebraizability is used for orthogonality only}
\label{warn:ch29-algebraizability-scope}
In \cite{BFJ15} the hypothesis \textup{(S4)} is invoked in exactly one
place, \cite[Appendix~A]{BFJ15}, and for exactly one purpose: to prove the
orthogonality property \eqref{eq:orthogonality} for an ample class.  It does
not enter the construction of the Monge--Amp\`ere operator
(\Cref{chap:bfj-ma}), the capacity theory (\Cref{chap:bfj-capacity}), the
energy calculus (\Cref{chap:bfj-energy}), the compactness or regularization
theorems (\Cref{chap:bfj-compactness}), the continuity of envelopes
(\Cref{chap:bfj-envelopes}), or the capacity and regularity arguments in the present chapter.
It enters only through the historical proof of orthogonality.  A
reader who wishes to check this may follow the ledger
\Cref{thm:orthogonality-provenance} and the table of
\Cref{sec:ch29-hypothesis-use}.
\end{warning}

\subsection{Removing algebraizability in residue characteristic zero}
\label{sec:ch29-remove-algebraizability}

\begin{theorem}\label{thm:bggjkm-remove-algebraizability}
Let $K$ be complete discretely valued with residue characteristic zero,
and let $X$ be smooth connected projective over $K$.  Every closed
model $(1,1)$-form
with ample de Rham class satisfies orthogonality.  Consequently
\Cref{cor:bfj-original-continuous-solution} holds without assuming
that $X$ is algebraizable over a global function field.
\end{theorem}

\begin{proof}[Imported intersection-theoretic input]
The differentiability theorem for non-Archimedean volumes identifies the
derivative at a continuous semipositive metric with integration against its
Chambert--Loir measure.  The relative volume of a metric and its semipositive
envelope vanishes.  Differentiating this equality along the segment toward
the obstacle gives
\[
 \int(f-P_\omega f)\MA_\omega(P_\omega f)=0.
\]
Rational approximation and homogeneity pass from line bundles to arbitrary
real closed forms with ample class.  This is
\cite[Theorem~6.3.3 and Corollary~6.3.4]{BGGJKM20}.
\end{proof}

This result removes one global-origin hypothesis.  It does not remove
discreteness, smoothness, residue characteristic zero, or support on a dual
complex from the cited continuous-solution theorem.

\begin{corollary}
\label{cor:ch29-bggjkm-cy}
Let $K$ be complete discretely valued with residue characteristic zero,
let $X$ be a smooth connected projective $K$-variety of dimension $n$,
let $\omega$ be a closed model $(1,1)$-form with ample de Rham class, and let
$\mu$ be a positive Radon measure of mass $V=(\{\omega\}^n)$ supported
on the dual complex of a projective SNC model.  Then there is an
$\omega$-psh function $\varphi\in C^0(X^{\An})$, unique modulo
constants, with $\MA_\omega(\varphi)=\mu$.  No algebraizability
hypothesis is made.
\end{corollary}

\begin{proof}
If $\omega$ is semipositive, combine
\Cref{thm:bggjkm-remove-algebraizability} with
\Cref{cor:bfj-original-continuous-solution}.  For a general closed
model form, choose a positive representative $\omega'$ of its ample
generic class on a higher model.  The exact sequence
\Cref{thm:ch20-ddc-exact} gives a real model function $g$ with
$\omega'=\omega+\ddc g$.  The equation transports by
\[
 \varphi=u+g,\qquad
 \omega+\ddc\varphi=\omega'+\ddc u .
\]
Apply the semipositive case to $\omega'$ and the same measure, enlarging
the SNC model if needed.  Both $g$ and $u$ are continuous, and uniqueness
transports by the same equality.  This reference change supplies the
semipositivity missing from a direct application of the preceding
corollary.
\end{proof}

\begin{remark}[The exact shape of the 2020 improvement]
\label{rem:ch29-2020-improvement}
In the 2015 proof the global-origin hypothesis \textup{(S4)} supplies
orthogonality, then differentiability of $E_\omega\circ P_\omega$
and variational existence.  The 2020 proof uses two separate inputs:
\[
 \begin{gathered}
 \text{continuous semipositive metric}
   \ \Longrightarrow\ \text{relative-volume derivative},\\
 \text{continuous envelope and preservation of small sections}\\
 \qquad+\ \text{that derivative}
   \ \Longrightarrow\ \text{orthogonality}.
 \end{gathered}
\]
The derivative is \cite[Theorem~5.4.3]{BGGJKM20} and does not follow
from continuity of all envelopes.  Preservation of small sections is
Proposition~6.2.1; applying the derivative at the continuous envelope
gives Theorem~6.3.2.  Rational approximation and homogeneity give the
real-class statement in Theorem~6.3.3.  Under \textup{(S1)}--\textup{(S3)},
BFJ supplies the required envelope continuity.
The continuous supported-solution theorem retains \textup{(S5)}.
In equal characteristic $p$, the cited test-ideal theorem supplies
envelope continuity under geometric-origin and resolution hypotheses;
its orthogonality theorem applies the same Theorem~6.3.2.
\end{remark}

\subsection{The support hypothesis \textup{(S5)}: what is known}
\label{sec:ch29-support-status}

The cited supported continuous-solution theorems retain
\textup{(S5)}.  It is a sufficient regularity condition, not a necessary
one: the bounded contrast in \Cref{ex:ch28-unbounded-dual-energy} is
continuous with a measure carried by no dual complex.
The following distinctions describe what is established here.

\emph{It cannot be dropped outright.}  Under the full standing hypotheses,
\Cref{ex:ch29-unbounded-not-supported} exhibits on $\mathbf P^1_K$ a
measure $\mu$ of mass $V$ whose unique finite-energy solution is
unbounded.  Its finite dual energy follows directly from the computation at its
finite-energy potential in \Cref{prop:ch28-dual-energy-of-ma}.  So ``$\mu$ has finite energy'' is
strictly weaker than ``$\mu$ has a continuous solution'', and no theorem
of the form ``finite energy $\Rightarrow$ continuity'' is available.

\emph{A capacity condition gives a separate boundedness result.}
\Cref{thm:ch29-kolodziej-linfty} shows that
$(A,\alpha)$-domination by capacity bounds the oscillation of the
solution.  That hypothesis is not implied by \textup{(S5)} in dimension
$\geq2$ as far as this book proves
(\Cref{rem:ch29-floor-open}), and it does not by itself yield continuity
(\Cref{rem:ch29-linfty-scope}).  The two displayed arguments prove different conclusions; no
general implication between their hypotheses is established here.

\emph{It is not a statement about retractions.}  See
\Cref{warn:ch29-support-not-skeletal} and
Item~\cref{itm:ch29-notsupp:factor} of \Cref{sec:ch29-support-not}.

\emph{On curves it may be weakened using Thuillier's theory.}  For
$\dim X=1$ a positive measure of the right mass always has a semipositive
solution, possibly singular, and the retraction of a semipositive metric to
a skeleton is again semipositive
(\Cref{rem:ch23-curve-specialness}); support on a skeleton then upgrades
the singular solution to a continuous one directly, without capacity.  This
is the route of \Cref{thm:curves-arbitrary-field-continuity}, and it is the
only place in this chapter where \textup{(S5)} is used in a way that has
nothing to do with Dini's theorem.

\section{Later continuity theorems and their extra hypotheses}
\label{sec:ch29-later}

The three theorems collected in this section extend
\Cref{cor:bfj-original-continuous-solution} in three different directions.
They are stated separately, each with its own date, source, and hypothesis
list, and none of them is a special case of another.  The first drops
essentially every hypothesis except the dimension; the second changes the
characteristic and pays with geometric origin and resolution; the third
changes it again and pays with two conditional resolution statements.  All
three retain \textup{(S5)}.

\subsection{Curves over arbitrary complete fields}
\label{sec:ch29-later-curves}

Curves are exceptional because their Monge--Amp\`ere operator is the graph
Laplacian.  No higher-dimensional resolution theorem is needed.

\begin{center}
\small
\begin{tabular}{p{.30\textwidth}|p{.62\textwidth}}
\toprule
\textbf{Dropped}&\textbf{Retained or added}\\
\midrule
discreteness of $v_K$
 &$\dim X=1$\\
residue characteristic zero
 &$K$ complete, nontrivially valued\\
higher-dimensional resolution
 &$L$ ample, $X$ geometrically connected; for the equation,
  $\mu$ supported on a strictly semistable skeleton\\
algebraizability \textup{(S4)}
 &mass $\mu(X^{\An})=\deg L$\\
\bottomrule
\end{tabular}
\end{center}

\begin{theorem}\label{thm:curves-arbitrary-field-continuity}
Let $K$ be any complete nontrivially valued non-Archimedean field, let
$X$ be a smooth projective geometrically connected curve, and let
$L$ be ample.  The envelope over semipositive rational model metrics
of every continuous metric on $L^{\An}$ is a
continuous semipositive metric.  If a positive measure of mass $\deg L$
is supported on the skeleton of a strictly semistable model, the associated
Monge--Amp\`ere equation has a continuous semipositive solution, unique up
to scale.
\end{theorem}
\begin{proof}[Source boundary]
The envelope assertion is \cite[Theorems~1.1 and~3.1]{GJKM19}; the latter
states that for a closed $(1,1)$-form of nef de Rham class on a smooth
projective curve, the envelope of a continuous obstacle is a uniform limit
of psh model functions.  In the notation of
\cite[Proposition~3.7]{GJKM19}, the curve-specific positivity test used in
that proof is
\begin{equation}
\label{eq:curve-slope-criterion}
 \sum_{\nu\in T_x(\Gamma)}
   w_x(\nu)\lambda_{x,\nu}(F)+\deg(\theta|_{C_x})\geq0
 \qquad(x\in\Gamma).
\end{equation}
For the equation, the deduction following
\cite[Proposition~1.2]{GJKM19} starts with Thuillier's possibly singular
semipositive solution, retracts it to the skeleton carrying $\mu$, and uses
the curve-specific facts that retraction preserves semipositivity and
pushes the curvature measure forward.  Those facts are proved in
\cite[Propositions~3.5 and~3.8]{GJKM19}.  The same discussion invokes the
non-Archimedean Calabi--Yau uniqueness theorem for uniqueness up to scale.
Thus both assertions are imported in full from the cited source; no
reconstruction of its descent, value-group approximation, or comparison
between Thuillier and model semipositivity is claimed here.
\end{proof}

\begin{remark}[Source, date, and the division of labour]
\label{rem:ch29-curve-source}
The envelope half of \Cref{thm:curves-arbitrary-field-continuity} is
\cite[Theorem~1.1]{GJKM19}, published in 2019, for an ample line bundle.
The rational-model envelope convention is explicit in
\Cref{thm:ch23-gjkm-curves}; the retraction statements used in its proof are
\cite[Proposition~1.2 and Propositions~3.5, 3.8]{GJKM19}.  The
Monge--Amp\`ere half is the deduction sketched in the introduction of
\cite{GJKM19}, resting on the existence of a possibly singular semipositive
solution in Thuillier's potential theory on a Berkovich curve, and on the
uniqueness statement quoted there.  The counterexample to the
higher-dimensional analogue of the retraction statement is due to Burgos Gil
and Sombra and appears in the appendix to \cite{GJKM19}.  Nothing in this
subsection uses \Cref{chap:bfj-capacity} or \Cref{chap:bfj-energy}: on a
curve the higher-dimensional capacity argument is replaced by Thuillier's
solution and the curve-specific retraction theorem.
\end{remark}

\begin{remark}[Where the ``nef'' and ``ample'' hypotheses sit]
\label{rem:ch29-curve-nef-ample}
Ampleness is retained in both assertions of
\Cref{thm:curves-arbitrary-field-continuity}.  Matching the total masses
is what makes $\mu-\mu_0$ have mass zero; positivity of $\deg L$
is a separate hypothesis.  Generic nefness alone is not used as a
model-existence assertion in degree zero.
\end{remark}

\subsection{Equicharacteristic positive characteristic}
\label{sec:ch29-later-positive}

The positive-characteristic theorem replaces multiplier ideals by test
ideals.  Its geometric-origin and resolution assumptions should be kept
visible.  We display them in the form in which they are stated in the
source, as the three assumptions \textup{(A1)}--\textup{(A3)} of
\cite[\S9]{GJKM19}:
\begin{equation}
\label{eq:ch29-gjkm-assumptions}
\begin{aligned}
 \textup{(A1)}\quad&\text{$X$ is a smooth projective $K$-variety of
   dimension $n$,}\\
 &\text{of geometric origin from a $d$-dimensional family}\\
 &\text{over a perfect field $k$ of characteristic $p>0$; }\\
 \textup{(A2)}\quad&\text{resolution of singularities holds over $k$ in
   dimension $d+n$;}\\
 \textup{(A3)}\quad&\text{embedded resolution of singularities holds over
   $k$ in dimension $d+n$.}
\end{aligned}
\end{equation}
Here $K$ is complete discretely valued of characteristic $p$---equal
characteristic, not merely residue characteristic $p$---and geometric
origin is \Cref{def:ch23-geometric-origin}.  Assumption \textup{(A2)} is
what \Cref{thm:ch23-gjkm-positive-characteristic} needs for continuity of
envelopes; \textup{(A3)} is added on top of it, and is used for one purpose
only, namely to make projective SNC models cofinal so that the skeletal
definition of psh functions and the whole calculus of
\Cref{chap:bfj-psh,chap:bfj-ma,chap:bfj-capacity,chap:bfj-energy} can be
reconstructed.

\begin{theorem}\label{thm:gjkm-positive-characteristic-cy}
Let $K$ be complete discretely valued of characteristic $p>0$.  Let
$X$ be a smooth projective variety of dimension $n$, of geometric origin
from a $d$-dimensional family over a perfect field $k$ of characteristic
$p$.  Assume resolution of singularities and embedded resolution of
singularities over $k$ in dimension $d+n$.  Let $L$ be ample and let
$\mu$ be a positive measure of mass $(L^n)$, supported on the skeleton
of a projective SNC model.  Then
\[
                    c_1(L,\|\cdot\|)^n=\mu
\]
has a continuous semipositive solution, unique up to scaling.
\end{theorem}
\begin{proof}
We record the source's deduction and its exact inputs.  Resolution
provides regular projective models for spread-out data.  For a model
line bundle $\mathcal L$ of the ample generic polarization, the base
ideals $\mathfrak a_m$ and asymptotic test ideals
$\mathfrak b_m=\tau(\mathfrak a_\bullet^m)$ satisfy containment,
subadditivity and uniform global generation.  These are the algebraic
inputs of \cite[Theorem~1.3 and Theorem~4.6]{GJKM19}.  Under the precise
spread-out hypotheses of its Assumption~7.1, the complete estimates in
Theorem~7.3 give
\[
       m^{-1}\log|\mathfrak a_m|\longrightarrow P_\theta(0)
       \quad\text{uniformly}.
\]
The normalized functions are monotone along divisibility, not asserted
monotone at all consecutive levels.  Model-reference changes,
homogeneity and rational approximation give continuous-obstacle envelopes
for real ample classes in \cite[Theorem~8.2]{GJKM19}.

Under \textup{(A1)}--\textup{(A3)}, projective SNC models are cofinal.
The source identifies its model-envelope convention with the BFJ
envelope; its Corollary~9.1 then supplies a decreasing net of psh model
regularizations.  A decreasing family of continuous obstacles and
independently chosen model approximants does not automatically supply
a decreasing diagonal net; the cited regularization theorem provides
that step.

Theorem~9.2 of the source first reduces to rational classes, then applies
\cite[Theorem~6.3.2]{BGGJKM20} at the continuous envelope.  This is the
orthogonality theorem; relative-volume differentiability is the separate
Theorem~5.4.3 of the latter source.  Orthogonality yields the projected
energy derivative needed to maximize
\[
                 E_\theta(u)-\int u\,\rd\mu
\]
on the normalized compact psh space.  The fixed skeletal support gives
continuous pairing; upper semicontinuity of energy gives a maximizer;
Euler--Lagrange identifies its measure.  The capacity and Dini argument
then proves continuity.  Uniqueness imports the BFJ15
Proposition~8.2/Yuan--Zhang Hodge argument, not the comparison principle.
This is \cite[Theorem~9.3 and its proof]{GJKM19}.
\end{proof}

Here ``of geometric origin'' includes the descent of the valued field and
variety needed to apply global test ideals.  Under resolution alone the
semipositive envelope is continuous; embedded resolution makes SNC models
cofinal and permits the full BFJ calculus.  Since resolution and embedded
resolution are known in dimension three over a perfect field, the theorem
is unconditional for a smooth projective surface of geometric origin from a
one-dimensional family.  These are Theorems~1.4, 1.5, and~9.3 of
\cite{GJKM19}.

\begin{corollary}
\label{cor:ch29-gjkm-surface-cy}
Let $K$ be complete discretely valued of characteristic $p>0$, and let
$X$ be a smooth projective \emph{surface} over $K$ of geometric origin
from a one-dimensional family over a perfect field $k$.  Then
\textup{(A2)} and \textup{(A3)} of \eqref{eq:ch29-gjkm-assumptions} hold
unconditionally, and consequently the conclusion of
\Cref{thm:gjkm-positive-characteristic-cy} holds with no resolution
hypothesis.
\end{corollary}

\begin{proof}[Source]
Here $n=2$ and $d=1$, so $d+n=3$, and both resolution and embedded
resolution of singularities in dimension three over a perfect field are the
theorems of Cossart and Piltant quoted as \cite[Theorem~6.3]{GJKM19}.  For
the envelope half this is \Cref{cor:ch23-gjkm-surfaces}; for the
Monge--Amp\`ere half it is the corresponding specialization of
\cite[Theorem~9.3]{GJKM19}.  Nothing else changes.
\end{proof}

\begin{warning}{Equal characteristic $p$ is not residue characteristic
$p$}
\label{warn:ch29-equal-vs-residue}
\Cref{thm:gjkm-positive-characteristic-cy} requires
$\operatorname{char}K=p$, not merely $\operatorname{char}\widetilde K=p$.
Test ideals are defined by Frobenius on $K$-schemes, and there is no
Frobenius in mixed characteristic; this is why $\mathbf Q_p$ is outside
the scope of the theorem, and why the mixed-characteristic case
\Cref{thm:fgk-mixed-characteristic-cy} required a different circle of
ideas in the 2025 theorem.  The same point is made for envelopes in
\Cref{rem:ch23-gjkm-hypotheses} and instantiated in
\Cref{ex:ch23-outside-scope}.
\end{warning}

\subsection{Mixed characteristic}
\label{sec:ch29-later-mixed}

\begin{theorem}\label{thm:fgk-mixed-characteristic-cy}
Let $K$ be complete discretely valued of mixed characteristic $(0,p)$,
and let $X$ be a smooth projective $K$-variety.  Assume:
\begin{enumerate}
\item resolution of singularities holds for projective
      $K^\circ$-models of $X$;
\item embedded resolution holds for regular projective
      $K^\circ$-models of $X$.
\end{enumerate}
If $\theta$ has ample de Rham class and $\mu$ has mass
$(\{\theta\}^n)$ and is supported on the skeleton of a projective SNC
model, then there is a continuous $\theta$-psh function $\varphi$,
unique modulo constants, satisfying
\[
                         \MA_\theta(\varphi)=\mu.
\]
\end{theorem}
\begin{proof}
The algebraic envelope input is
\cite[Theorem~5.1]{FGK25}.  It assumes a regular projective model on
which the ample generic polarization has an ample extension
$\mathcal A$.  For another extension $\mathcal L=\mathcal O(D)$
of that polarization, let $\mathfrak a_m$ be its base ideal and put
\[
 \mathfrak b_m=\tau(\mathcal O_{\mathcal X},m\|D\|),\qquad
 \mathfrak a'_m=\mathfrak a_m\,\tau(\mathcal O_{\mathcal X},0).
\]
The fixed correction factor $\tau(\mathcal O_{\mathcal X},0)$ cannot
be omitted.  The source proves, at sufficiently large divisible levels,
\[
 \mathfrak a'_m\subseteq\mathfrak b_m,\qquad
 \mathfrak b_{m\ell}\subseteq\mathfrak b_m^\ell,\qquad
 \mathfrak b_m\otimes\mathcal A^{m_0}\otimes\mathcal L^m
       \text{ globally generated},
\]
with $m_0$ independent of $m$.  The correction changes the normalized
model function by the fixed bounded model function
$m^{-1}\log|\tau(\mathcal O_{\mathcal X},0)|$.  The full estimates in
the cited proof therefore still give uniform convergence of
$m^{-1}\log|\mathfrak a_m|$ to $P_\theta(0)$.
The graded base-ideal relation gives superadditivity before division
by the index and monotonicity of the normalized functions along
divisibility; consecutive normalized terms need not be monotone.
This paragraph records the algebraic inputs and their correction,
rather than reproving the entire ideal-theoretic estimate.

For a continuous obstacle, uniform model approximation reduces to a
rational model obstacle.  Homogeneity clears its denominators.
The resolution and ample-extension lemma
\cite[Lemma~6.4]{FGK25} supplies the required regular projective model,
after a positive tensor power; the model-reference identity
$P_\theta(u)=P_{\theta+\ddc u}(0)+u$ then applies Theorem~5.1.
This is precisely \cite[Theorem~7.1 and its proof]{FGK25}.
With embedded resolution, SNC models are cofinal.  The source's
Section~8 explains the transfer of the BFJ psh, regularization and
Monge--Amp\`ere calculus and the real ample-class passage.
Continuous envelopes and
\cite[Theorems~6.3.2--6.3.3]{BGGJKM20} give orthogonality, with the
relative-volume derivative supplied separately by Theorem~5.4.3.

The variational functional attains a maximum because skeletal pairing
is continuous on normalized psh functions.  The projected derivative
gives the prescribed measure, and the supported capacity--Dini proof
gives continuity.  Uniqueness imports the BFJ15
Proposition~8.2/Yuan--Zhang Hodge argument.  This is
\cite[Theorem~8.1 and its full proof]{FGK25}; comparison alone is not
its uniqueness input.
\end{proof}

The proof uses perturbation-friendly test ideals in mixed characteristic to
prove continuity of envelopes, then invokes the BFJ variational and capacity
arguments.  This is \cite[Theorems~7.1 and~8.1]{FGK25}.  The result is
conditional on the two stated resolution properties.  Existing
three-dimensional resolution results do not automatically supply the
required projective morphisms, so one should not silently advertise the
mixed-characteristic surface case as unconditional.

\begin{remark}[The two resolution hypotheses are about $K^\circ$-models]
\label{rem:ch29-fgk-resolution}
The hypotheses of \Cref{thm:fgk-mixed-characteristic-cy} differ in kind
from \textup{(A2)} and \textup{(A3)} of
\eqref{eq:ch29-gjkm-assumptions}.  There, resolution is required over a
perfect field $k$ in a dimension $d+n$ determined by the family, and it
is a statement about varieties.  Here, resolution is required for
projective $K^\circ$-\emph{models}, that is, for schemes flat and
projective over a mixed-characteristic discrete valuation ring, and it is a
statement about arithmetic schemes.  The two are not comparable: the known
three-dimensional resolution results in mixed characteristic produce
regular models by proper modifications, but projectivity is not part of
their stated conclusion, and
the BFJ calculus of \Cref{chap:bfj-model-forms} is set up for projective
models throughout, since relative ampleness is used in
\Cref{prop:ch20-positive-representative} and hence in
\Cref{lem:divisorial-positive-capacity}.  This is the precise content of
the warning in the last paragraph above, and the reason why the
mixed-characteristic surface case is not advertised here as unconditional.
\end{remark}

\begin{remark}[Dates, in order]
\label{rem:ch29-dates}
For the record: the original theorem is \cite{BFJ15}, published in 2015
with the hypothesis list \eqref{eq:ch29-signature}; the curve theorem and
the equal-characteristic test-ideal theorem are \cite{GJKM19}, published in
2019; the removal of algebraizability is \cite{BGGJKM20}, published in
2020; the conditional mixed-characteristic theorem is \cite{FGK25},
published in 2025.  Each is stated above as a separate theorem with its own
hypotheses, and no statement in this chapter combines the hypothesis lists
of two of them.
\end{remark}

\subsection{Hypothesis atlas}
\label{sec:ch29-atlas}

\begin{description}[style=nextline,leftmargin=0pt,labelindent=0pt,
  labelsep=0pt,itemsep=.8\baselineskip,font=\normalfont\bfseries]
\item[BFJ 2015]
\emph{Input.} Multiplier ideals and algebraic orthogonality.
\par\emph{Remaining hypotheses.} Complete discrete valuation, residue
characteristic zero, smooth projective and algebraizable variety, ample
class, and a measure on one SNC dual complex.

\item[BGGJKM 2020]
\emph{Input.} Volume differentiability proves orthogonality.
\par\emph{Remaining hypotheses.} The same BFJ analytic setting, but without
algebraizability; dual-complex support remains.

\item[Curves]
\emph{Input.} Thuillier theory and graph retraction.
\par\emph{Remaining hypotheses.} The curve is smooth projective and
geometrically connected, the field is complete and nontrivially valued,
and the line bundle is ample.  For the equation the measure has mass
$\deg L$ and is carried by a semistable skeleton.

\item[GJKM 2019]
\emph{Input.} Equal-characteristic test ideals.
\par\emph{Remaining hypotheses.} Characteristic $p>0$, geometric origin,
resolution and embedded resolution in dimension $d+n$, and dual-complex
support.

\item[FGK 2025]
\emph{Input.} Mixed-characteristic test ideals.
\par\emph{Remaining hypotheses.} Mixed characteristic, resolution for
projective models, embedded resolution for regular projective models, and
dual-complex support.

\item[Toric or invariant setting]
\emph{Input.} Real Alexandrov theory.
\par\emph{Remaining hypotheses.} The invariant category and the precise
tropical hypotheses of \Cref{chap:cld-examples}.
\end{description}

The ledger separates continuity of envelopes from continuity of solutions.
An envelope theorem makes the variational calculus available.  Continuity of
the resulting solution still uses the location of the prescribed measure,
unless another criterion proves continuity.  An oscillation estimate
alone establishes boundedness, a weaker conclusion.

\section{What finite energy does not imply}

\paragraph{Finite energy versus regularity.}
Under the corresponding orthogonality hypotheses, finite energy implies
that $\mu$ is a strong limit of model Monge--Amp\`ere measures.  If,
more specifically, $\mu$ is supported on one SNC dual complex, then it
has a finite-energy solution and that solution is continuous.  \Cref{thm:ch28-finite-energy-existence} supplies an
$\mathcal E^1$-potential for every finite-energy measure in the standing
BFJ framework under orthogonality, by quantitative estimates and
compact energy-level pairing.  Completion alone is not that argument.
Skeletal support remains sufficient for continuity and is not necessary,
as the bounded contrast in \Cref{ex:ch28-unbounded-dual-energy} shows.

Three tempting extrapolations are therefore invalid:
\begin{enumerate}
\item\label{itm:ch29-invalid:polar}
      a measure that does not charge pluripolar sets need not have finite
      energy;
\item\label{itm:ch29-invalid:bounded}
      a finite-energy measure need not have a bounded or continuous
      potential;
\item\label{itm:ch29-invalid:weak}
      weak convergence of finite-energy measures does not by itself imply
      uniform convergence of their potentials.
\end{enumerate}

Assertion \cref{itm:ch29-invalid:bounded} is
\Cref{ex:ch29-unbounded-not-supported}, where the measure
\eqref{eq:ch29-boundary-measure} is a countable sum of atoms at divisorial
points---so it charges no pluripolar set whatsoever, by
\Cref{prop:ch29-skeleton-not-polar} applied on each
$\Delta_{\mathcal X_m}$---and yet its unique solution is unbounded.
Assertion \cref{itm:ch29-invalid:weak} is visible in the same example: the
truncations $\Phi_m$ of \Cref{ex:ch26-unbounded-finite-energy} are model
functions whose Monge--Amp\`ere measures converge weakly to $\mu$, while
$\Phi_m\to\varphi$ only pointwise, the convergence being uniform on no
neighbourhood of the point $T=0$.  The corresponding positive statement,
uniform convergence \emph{on a fixed dual complex} for a weakly convergent
family carried by that complex, is
\Cref{prop:fixed-skeleton-stability}.  No assertion of an exhaustive
list of stability results is made.

A fourth extrapolation, less often made but equally invalid, is
that continuity implies a model solution.  The quadratic solution in
\Cref{ex:ch29-curve-solution} is Lipschitz on its skeleton but not
piecewise affine; the atomic solution in \Cref{ex:ch29-curve-dirac}
is a model function.  Thus these examples separate model regularity,
without excluding Lipschitz regularity.

Assertion \cref{itm:ch29-invalid:polar} has a different elementary
example.  On $\mathbf P^1$ with the standard reference and mass one,
let $y_k$ be the divisorial point at depth $2^k$ on $R_0$ and put
\[
              \nu=\sum_{k\geq1}\frac1{k(k+1)}\delta_{y_k}.
\]
Every atom is divisorial, so $\nu$ charges no $\sigma$-pluripolar set.
For $N\geq1$, truncate the ray at $T_N=2^N$, moving all farther
atoms to $y_N$, and let $\nu_N$ be this pushforward.  Its graph
potential $v_N$, with $v_N(\eta)=0$, has slope minus the tail mass:
it is $-1$ on $[0,2]$, $-1/k$ on
$[2^{k-1},2^k]$ for $2\leq k\leq N$, and zero beyond $T_N$.
Extend constantly on attached branches.  These rational piecewise
affine slopes and breakpoints give an $\omega$-psh model function
with $\MA_\omega(v_N)=\nu_N$.  Since $v_N$ is constant beyond $T_N$,
its pairings with $\nu$ and $\nu_N$ agree.  The curve energy formula
and integration by parts give
\[
 \begin{aligned}
 D_N&:=\int_0^{T_N}|v_N'|^2\,\rd t
       =2+\sum_{k=2}^N\frac{2^{k-1}}{k^2}\longrightarrow+\infty,\\
 E_\omega(v_N)&=-\tfrac12D_N,\qquad
 \int v_N\,\rd\nu=\int v_N\,\rd\nu_N=-D_N .
 \end{aligned}
\]
Consequently the model-test definition gives
$E_\omega^\vee(\nu)\geq F_\nu(v_N)=D_N/2$ for every $N$,
hence $E_\omega^\vee(\nu)=+\infty$.  Avoiding polar mass alone
therefore does not imply finite dual energy.  This calculation is an
application of the established curve formulas, not a new research claim.

\begin{example}[Why Dini is localized]
\label{ex:ch29-dini-localized}
Let $\varphi_j\searrow\varphi$ be the BFJ regularization of a solution.
On every fixed dual complex the convergence is uniform because the limit is
continuous there.  Without a fixed complex carrying $\mu$, this skeletal argument
does not supply a single compact set on which uniform convergence gives
the almost-everywhere inequality required by
\Cref{lem:bfj-supported-domination}.  Choosing a different skeleton for
each index supplies only the bounds in \Cref{prop:ch29-one-complex}.
For some nonskeletal measures the solution is continuous by other means;
no impossibility of uniform approximation is asserted for every such measure.
\end{example}

\section{Worked examples}
\label{sec:ch29-worked}

The examples are organized in three groups.  On a curve segment we
compute a continuous quadratic solution and an atomic model solution,
including energies and capacity bounds.  In arbitrary dimension we apply
the capacity floor to a finite divisorial support.  Finally, the unbounded
ray example separates finite energy from boundedness and shows failure at
the borderline capacity exponent.

\subsection{A continuous solution on a curve skeleton}
\label{sec:ch29-ex-curve}

\begin{example}[The uniform measure on a segment of length four]
\label{ex:ch29-curve-solution}
Use the split curve setting above.  Start with a semistable model
with two components meeting at a split node $xy=\pi^4$; that model is
not regular.  Its regular resolution is a chain of four unit edges,
with three intermediate components, and retains the modulus segment
$\Gamma=[0,4]$ with endpoints $x_{E_0}$ and $x_{E_1}$.
This is the setup of \Cref{ex:ch25-segment-capacities}.
Let the generic line bundle $L$ have degree two, and let $\omega$
be the pullback of a model reference with endpoint degrees one and
degree zero on the intermediate components, so that
\[
 \mu_0=\MA_\omega(0)=\delta_{x_0}+\delta_{x_4},
 \qquad m_0=m_1=1,\qquad V=\deg L=2,\qquad n=1 .
\]
Let $\mu$ be the measure on $\Gamma$ with density $\tfrac12$ with
respect to arclength:
\[
 \rd\mu=\tfrac12\,\rd s\ \text{ on }[0,4],
 \qquad \mu(C^{\An})=\tfrac12\cdot4=2=V .
\]
Thus $\mu\in\mathcal M_V(\Gamma)$ and \Cref{cor:bfj-original-continuous-solution}
applies.

\emph{The solution.}  Write $\varphi=u\circ p$ with
$u:[0,4]\to\mathbf R$.  By \eqref{eq:ch25-segment-measure},
\[
 \MA_\omega(u\circ p)
 =\bigl(1+u'(0^+)\bigr)\delta_{x_0}
  +u''|_{(0,4)}
  +\bigl(1-u'(4^-)\bigr)\delta_{x_4} .
\]
Matching this with $\mu$ forces
\[
 u'(0^+)=-1,\qquad u''=\tfrac12\ \text{ on }(0,4),
 \qquad u'(4^-)=1,
\]
which is consistent, since $u'(4^-)-u'(0^+)=\tfrac12\cdot4=2$.  Hence
$u'(x)=-1+\tfrac x2$ and
\begin{equation}
\label{eq:ch29-curve-solution}
 u(x)=\frac{x^2}{4}-x+c .
\end{equation}
The function $u$ is convex, so its maximum on $[0,4]$ is attained at an
endpoint, and $u(0)=u(4)=c$; the normalization $\sup\varphi=0$ therefore
gives $c=0$.  Its minimum is at $x=2$, where $u(2)=1-2=-1$.  Thus
\begin{equation}
\label{eq:ch29-curve-solution-values}
 \varphi=\Bigl(\frac{x^2}{4}-x\Bigr)\circ p,
 \qquad \sup_{C^{\An}}\varphi=0,
 \qquad \inf_{C^{\An}}\varphi=-1,
 \qquad \operatorname{osc}\varphi=1 .
\end{equation}

\emph{Verification, by explicit model approximation.}  The function $u$ is
not piecewise affine, so $\varphi$ is not a model function, and the
computation above has to be justified by an approximation which we now carry
out with numbers.  For $N\geq1$ let $u_N$ be the piecewise affine
interpolant of $u$ at the $N+1$ points $4k/N$, $0\leq k\leq N$.
Since $u$ is convex, $u_N\geq u$, refining the partition decreases the
interpolant, so $u_{2N}\leq u_N$, and
\begin{equation}
\label{eq:ch29-curve-interpolation-rate}
 0\ \leq\ u_N-u\ \leq\ \frac18\,\sup|u''|\,\Bigl(\frac4N\Bigr)^2
 =\frac18\cdot\frac12\cdot\frac{16}{N^2}=\frac1{N^2}.
\end{equation}
The slope of $u_N$ on $[4k/N,4(k+1)/N]$ is the mean value
$u'(4k/N+2/N)=2k/N+1/N-1$, a rational number, and the breakpoints are
rational; so $u_N\circ p$ is a model function by
\Cref{lem:model-function-affine} and
\Cref{thm:skeleton-subdivision}.  Its Monge--Amp\`ere measure is computed by
\eqref{eq:ch25-segment-measure}: the outgoing slope at $0$ is
$1/N-1$, the outgoing slope at $4$ is $1-1/N$, and the slope jump at
each interior breakpoint is $2/N$, so
\begin{equation}
\label{eq:ch29-curve-approx-measure}
 \MA_\omega(u_N\circ p)
 =\frac1N\delta_{x_0}
  +\sum_{k=1}^{N-1}\frac2N\,\delta_{x_{4k/N}}
  +\frac1N\delta_{x_4},
\end{equation}
of total mass $\tfrac1N+(N-1)\tfrac2N+\tfrac1N=2=V$, as
\eqref{eq:ch24-mass} requires.  All the coefficients are nonnegative, so
each $u_N\circ p$ is $\omega$-psh.

Now $u_N\circ p$ decreases to $\varphi$ along the subsequence
$N=2^j$, so $\varphi\in\PSH(X,\omega)$ by
\Cref{thm:ch21-decreasing-stability}, and the convergence is uniform by
\eqref{eq:ch29-curve-interpolation-rate}; thus $\varphi$ is continuous and
is a uniform limit of $\omega$-psh model functions, as
\Cref{cor:ch29-uniform-model-limit} predicts.  Finally
$\MA_\omega(u_N\circ p)\to\MA_\omega(\varphi)$ weakly by
\Cref{prop:ch24-continuous-uniform}\,\cref{itm:ch24-cu:unif}, while the
atomic measures \eqref{eq:ch29-curve-approx-measure}, whose atoms sit at
spacing $4/N$ with common mass $2/N=\tfrac12\cdot\tfrac4N$, converge
weakly to the measure of density $\tfrac12$ on $[0,4]$.  Hence
$\MA_\omega(\varphi)=\mu$, as claimed.

\emph{Energy.}  By \eqref{eq:ch26-graph-energy},
\[
 E_\omega(\varphi)
 =\int_\Gamma u\,\rd\mu_0-\frac12\int_0^4|u'|^2\,\rd s
 =\bigl(u(0)+u(4)\bigr)
  -\frac12\int_0^4\Bigl(\frac x2-1\Bigr)^2\rd s .
\]
The first term is $0$.  For the second, substituting $y=x/2-1$ with
$\rd s=2\,\rd y$ and $y\in[-1,1]$,
\[
 \int_0^4\Bigl(\frac x2-1\Bigr)^2\rd s=2\int_{-1}^1y^2\,\rd y=\frac43 ,
\]
so
\begin{equation}
\label{eq:ch29-curve-energy}
 E_\omega(\varphi)=-\frac23 .
\end{equation}
In particular $\varphi\in\cE^1(C,\omega)$, as it must be, being bounded.

\emph{The Ko\l odziej bound, evaluated.}  We check
\Cref{cor:ch29-kolodziej-floor} on these numbers.  By
\eqref{eq:ch29-curve-point-capacity}, the capacity of the point $x_c$ is
$\min\{1,1/c\}+\min\{1,1/(4-c)\}$, and the exact infimum of this quantity
over $c\in[0,4]$ is $1$, attained at the midpoint; the three-case
verification is written out in \Cref{rem:ch29-floor-not-sharp}.  So the
capacity floor \eqref{eq:ch29-capacity-floor} holds with $c=1$, and
\Cref{lem:ch29-capacity-floor} gives $(A,\alpha)$-domination with
$A=Vc^{-(1+\alpha)}=V=2$ for every $\alpha>0$.  The constant
$C_\omega$ of
\eqref{eq:ch25-normalized-bound} may be taken to be $4$ here: for
$v\in\PSH(C,\omega)$ with $\sup v=0$, the restriction $v|_\Gamma$ is
convex with $v'(0^+)\geq-1$ and $v'(4^-)\leq1$, the supremum is attained
at an endpoint, and therefore
\[
 \int(-v)\,\omega=-v(x_0)-v(x_4)\leq4 ,
\]
the extreme case being $v(x)=-x$ or $v(x)=x-4$.  Hence
\begin{equation}
\label{eq:ch29-curve-kolodziej-numbers}
 \operatorname{osc}\varphi
 \ \leq\ (C_\omega+nV)\,A^{1/\alpha}2^{\,1/\alpha^2}+2
 \ =\ 6\cdot2^{1/\alpha}2^{1/\alpha^2}+2
 \ \xrightarrow[\ \alpha\to\infty\ ]{}\ 8 .
\end{equation}
So the estimate gives $\operatorname{osc}\varphi\leq8$, while the true
value is $1$.  The factor of eight is the price of the three sources of
slack listed in \Cref{rem:ch29-linfty-slack}; the estimate is nonetheless
of the right form, since it depends on $\mu$ only through $V$ and the
capacity floor.
\end{example}

\begin{example}[An atomic measure on the same skeleton]
\label{ex:ch29-curve-dirac}
Keep the data of \Cref{ex:ch29-curve-solution} and replace $\mu$ by
$\mu'=2\delta_{x_2}$, the Dirac mass of total mass $V=2$ at the midpoint
of $\Gamma$.  Matching \eqref{eq:ch25-segment-measure} again forces
\[
 u'(0^+)=-1,\qquad u''=2\delta_2,\qquad u'(4^-)=1 ,
\]
so $u'\equiv-1$ on $(0,2)$ and $u'\equiv1$ on $(2,4)$; with
$\sup u=0$,
\begin{equation}
\label{eq:ch29-curve-dirac-solution}
 u(x)=\begin{cases}-x,&0\leq x\leq2,\\ x-4,&2\leq x\leq4,\end{cases}
 \qquad
 \varphi'=u\circ p,
 \qquad \inf\varphi'=u(2)=-2 .
\end{equation}
This time $u$ is piecewise affine with integral slopes and integral
breakpoint, so $\varphi'$ \emph{is} an $\omega$-psh model function, by
\Cref{lem:model-function-affine} and the criterion recalled in the proof of
\Cref{cor:ch25-curve-capacity-formula}: the midpoint $x_2$ is already a vertex of the regular resolution,
and this function is determined on that model.  Its energy is
\[
 E_\omega(\varphi')
 =\bigl(u(0)+u(4)\bigr)-\frac12\int_0^4|u'|^2\rd s
 =0-\frac12\cdot4=-2 .
\]
The Ko\l odziej bound is the same as in
\eqref{eq:ch29-curve-kolodziej-numbers}, namely $8$, and the true
oscillation is $2$.

Two features are worth recording.  First, the solution here is a model
function, whereas in \Cref{ex:ch29-curve-solution} it is not; support on a
dual complex does not distinguish the two cases, and no theorem of this
chapter predicts which occurs.  Second, both solutions factor through the
retraction $p$, which is a phenomenon of dimension one only:
\Cref{warn:ch29-support-not-skeletal} records that in higher dimension a
solution of a skeleton-supported equation need not be invariant under any
SNC retraction.
\end{example}

\subsection{A measure carried by finitely many divisorial points}
\label{sec:ch29-ex-divisorial}

\begin{example}[A Dirac mass in arbitrary dimension]
\label{ex:ch29-divisorial-dirac}
Let $X$ and $\omega$ satisfy the standing hypotheses, let
$x=x_E$ be a divisorial point, and let $\mu=V\delta_x$, a measure of mass
$V$ supported on the dual complex of any SNC model having $E$ as a
component of its special fibre.  Then
\Cref{cor:bfj-original-continuous-solution} produces a continuous
$\omega$-psh $\varphi_x$ with
$\MA_\omega(\varphi_x)=V\delta_x$, normalized by $\sup\varphi_x=0$; it
differs by a constant from the Green function in
\eqref{eq:green-envelope-divisorial}.  That pointwise supremum is
proved continuous in \cite[Lemma~8.5 and equation~(8.7)]{BFJ15}
and is normalized to have value zero at $x$, whereas the solution
here is normalized by supremum zero.

The Ko\l odziej estimate applies with explicit data.  Put
$c:=\Capa_\omega(\{x\})$, which is positive by
\Cref{lem:divisorial-positive-capacity}.  Choose on a dominating model the
admissible potential $\psi$ in that lemma's proof, whose curvature has a
relatively ample representative.  The intersection number
$\lambda:=\bigl((\omega+\ddc\psi)|_E\bigr)^n>0$ in
\eqref{eq:ch25-divisorial-atom} satisfies
$c\geq b_E\lambda$ with $b_E$ the multiplicity of $E$ in the special
fibre.  Since $\Supp\mu=\{x\}$, \Cref{lem:ch29-capacity-floor} gives
$(A,\alpha)$-domination with $A=Vc^{-(1+\alpha)}$ for every
$\alpha>0$, and \eqref{eq:ch29-kolodziej-bound} reads
\begin{equation}
\label{eq:ch29-dirac-bound}
 \operatorname{osc}_{X^{\An}}\varphi_x
 \ \leq\ (C_\omega+nV)\,
   \bigl(Vc^{-(1+\alpha)}\bigr)^{1/\alpha}2^{\,n/\alpha^2}+2 .
\end{equation}
For every $c>0$, letting $\alpha\to\infty$ gives
$\operatorname{osc}\varphi_x\leq(C_\omega+nV)/c+2$, by
\eqref{eq:ch29-kolodziej-floor-bound}.  In fact $c\leq V$ because
every admissible measure has mass $V$, and the variable factor is
\[
 c^{-1}\exp\left(\frac{\log(V/c)}{\alpha}
                 +\frac{n\log2}{\alpha^2}\right).
\]
For $n\geq1$ it decreases toward $1/c$ as $\alpha$ increases;
there is no finite minimizing exponent when $c<1$.
The bound records a dependence on the capacity floor, not a proof that
its numerical growth is sharp.

The same applies verbatim to any measure
$\mu=\sum_{i=1}^Nc_i\delta_{x_i}$ with $x_i$ divisorial,
$c_i\geq0$ and $\sum_ic_i=V$: take
$c=\min_i\Capa_\omega(\{x_i\})>0$.  This supplies a capacity-floor application in every dimension of
the standing framework, without a separate uniform-floor assumption.
\end{example}

\subsection{What the support hypothesis buys: a measure outside its scope}
\label{sec:ch29-ex-boundary}

The next example is the boundary marker of the chapter.  It shows that under
the \emph{full} standing hypotheses---complete discretely valued field,
residue characteristic zero, smooth projective $X$, ample class, and a
measure of the correct total mass which moreover has finite energy---the
solution can fail to be bounded, hence fail to be continuous.  The measure violates \textup{(S5)} and the positive-power capacity
condition.  Neither finite energy nor the mass constraint supplies
continuous regularity.

\begin{example}[A finite-energy measure with an unbounded solution]
\label{ex:ch29-unbounded-not-supported}
Work on $X=\mathbf P^1_K$ with $L=\cO(1)$ and $\omega$ the curvature
of the standard model metric, so that $\mu_0=\MA_\omega(0)=\delta_\eta$ at
the Gauss point and $V=1$, $n=1$.  Let
$\mathbf t=-\varphi_0:X^{\An}\to[0,+\infty]$ be the coordinate of
\Cref{ex:ch25-pluripolar-computed}, so that $\{\mathbf t\geq s\}$ is the
closed ball $\{|T|\leq\rme^{-s}\}$, and let $\varphi=u(\mathbf t)$ be
the potential constructed in \Cref{ex:ch26-unbounded-finite-energy}, where
$u$ is piecewise affine with $u(0)=0$ and slope $-1/k$ on
$[k-1,k]$.  Recall from there that
\begin{equation}
\label{eq:ch29-boundary-measure}
 \mu:=\MA_\omega(\varphi)
 =\sum_{k\geq1}\frac1{k(k+1)}\,\delta_{\mathbf t=k},
 \qquad \mu(X^{\An})=1=V,
\end{equation}
where each atom denotes the canonical ray point at depth $k$.
Also $\varphi\in\cE^1(X,\omega)$ and
$\varphi(\mathbf t=k)=-H_k\to-\infty$, so $\varphi$ is unbounded and in
particular not continuous.  By \Cref{thm:ch29-uniqueness} it is the unique
solution modulo constants; there is therefore no continuous solution of
$\MA_\omega(\cdot)=\mu$.

\emph{$\mu$ is supported on no dual complex.}  The support of $\mu$ is
the closure of $\{\mathbf t=k:k\geq1\}$, which is that set together with
the rational point $T=0$.  A dual complex $\Delta_{\mathcal X}$ is a
compact polyhedral subset of the quasimonomial locus, and $T=0$ is a point
of type one, hence lies in no dual complex; moreover
\Cref{prop:ch29-skeleton-not-polar} shows directly that
$T=0$, being the polar set of $\varphi_0$, lies in no dual complex.
Since $\mu(\{T=0\})=0$ this alone is not decisive, so we argue with the
atoms.  Let $\mathcal X$ be any SNC model of $\mathbf P^1_K$.  Its dual
complex $\Delta_{\mathcal X}$ is compact and contained in the
quasimonomial locus, so
$\Delta_{\mathcal X}\cap R_0$ is a compact subset of
$R_0\setminus\{T=0\}$, which the parameter $\mathbf t$ identifies with
$[0,+\infty)$; hence there is an integer $m=m(\mathcal X)\geq1$ with
\[
 \Delta_{\mathcal X}\cap R_0\subseteq\{\mathbf t\leq m\} .
\]
For the standard tower of blowups one has exactly
$\Delta_{\mathcal X_m}=\{0\leq\mathbf t\leq m\}$.  A measure carried by
$\Delta_{\mathcal X}$ therefore gives no mass to the atoms of
\eqref{eq:ch29-boundary-measure} lying at $\mathbf t=k>m$; but
\[
 \mu\bigl(\{\mathbf t>m\}\bigr)=\sum_{k>m}\frac1{k(k+1)}
 =\frac1{m+1}>0
\]
for every $m$.  Hence $\mu$ is carried by no single dual complex, and
\Cref{prop:ch29-one-complex} applies to the tower
$(\mathcal X_m)$ with $\eta_m=1/(m+1)$.

\emph{$\mu$ fails the capacity-domination hypothesis for every
$\alpha>0$.}  Take $E_m=\{\mathbf t\geq m\}$.  By
\eqref{eq:ch25-p1-capacity},
\[
 \Capa_\omega(E_m)=\min\Bigl\{1,\frac1m\Bigr\}=\frac1m
 \quad(m\geq1),
 \qquad
 \mu(E_m)=\sum_{k\geq m}\frac1{k(k+1)}=\frac1m .
\]
Hence for every $A>0$ and $\alpha>0$,
\begin{equation}
\label{eq:ch29-borderline-ratio}
 \frac{\mu(E_m)}{\Capa_\omega(E_m)^{1+\alpha}}
 =\frac{1/m}{(1/m)^{1+\alpha}}=m^{\alpha}
 \ \xrightarrow[\ m\to\infty\ ]{}\ +\infty ,
\end{equation}
so \eqref{eq:ch29-capacity-domination} fails.  This is consistent with
\Cref{thm:ch29-kolodziej-linfty}, which would otherwise force $\varphi$ to
be bounded.
\end{example}

\begin{example}[The exponent $\alpha=0$ is not allowed]
\label{ex:ch29-borderline}
The computation \eqref{eq:ch29-borderline-ratio} shows more: in
\Cref{ex:ch29-unbounded-not-supported} one has \emph{exactly}
\[
 \mu(E_m)=\Capa_\omega(E_m)\qquad\text{for every }m\geq1 ,
\]
and in fact the same domination holds for every Borel set.  Write $x_k$ for the canonical point at depth $k$ on the ray
$R_0$, rather than for the full level set $\{\mathbf t=k\}$.  If $\mu(E)>0$, let
$m=\min\{k\geq1:x_k\in E\}$.  Then
\[
 \mu(E)\leq\sum_{k\geq m}\frac1{k(k+1)}=\frac1m.
\]
On the other hand the admissible potential
$\psi_m=-\min\{\mathbf t/m,1\}$ satisfies
\[
 \MA_\omega(\psi_m)
  =\Bigl(1-\frac1m\Bigr)\delta_\eta+\frac1m\delta_{x_m},
\]
by the one-dimensional slope formula.  Hence monotonicity of capacity gives
\[
 \Capa_\omega(E)\geq\Capa_\omega(\{x_m\})\geq\frac1m\geq\mu(E).
\]
The assertion is trivial when $\mu(E)=0$.  Thus $\mu$ satisfies the
borderline domination
$\mu(E)\leq A\,\Capa_\omega(E)^{1+\alpha}$ globally, with $A=1$ and
$\alpha=0$, while its potential is unbounded.
Therefore no version of \Cref{thm:ch29-kolodziej-linfty} with
$\alpha=0$ can hold.  This establishes the need for a stronger
hypothesis than this borderline one, not optimality of the threshold
or its rate in \Cref{rem:ch29-iteration-constants}.

The same example calibrates the iteration lemma.  With
$f(s)=\Capa_\omega(\{\varphi<-s\})$, the identity
$\varphi(\mathbf t=k)=-H_k$ gives
$\{\varphi<-H_k\}=\{\mathbf t>k\}$, so
$f(H_k)=1/k$ by \eqref{eq:ch25-p1-capacity-open}; since
$H_k\sim\log k$, the profile decays like
$f(s)\approx\rme^{-s}$, which is compatible with
\eqref{eq:ch29-profile-decay} and never reaches zero.  An iteration obeying
\eqref{eq:ch29-iteration-hypothesis} with $\alpha>0$ would force
$f$ to vanish identically beyond a finite level; the example shows that
\eqref{eq:ch29-iteration-hypothesis} itself must fail, and
\eqref{eq:ch29-borderline-ratio} says where.
\end{example}

\begin{remark}[What this example does and does not show]
\label{rem:ch29-boundary-scope}
\Cref{ex:ch29-unbounded-not-supported} does not contradict any theorem in
this chapter, and it is not a counterexample to a conjecture.  It shows
three things.  First, hypothesis \textup{(S5)} in
\Cref{thm:bfj-supported-regularity} cannot be replaced by ``$\mu$ has
finite energy'', nor by ``$\mu$ charges no pluripolar set'': the measure
\eqref{eq:ch29-boundary-measure} is atomic on divisorial points and charges
no pluripolar set at all.  Second, it cannot be replaced by ``$\mu$ is
carried by the quasimonomial locus'', since every atom of
\eqref{eq:ch29-boundary-measure} is a divisorial point.  Third, it cannot be
replaced by ``$\mu$ is carried by a countable union of dual complexes'',
since $\mu$ is carried by $\bigcup_m\Delta_{\mathcal X_m}$, nor by any
compactness requirement on $\Supp\mu$, which is automatic on the compact
space $X^{\An}$.  What fails is precisely carriage by \emph{one} dual
complex, and \Cref{prop:ch29-one-complex} quantifies the deficit as
$\eta_m=1/(m+1)$.
\end{remark}

\section{Scope: what each hypothesis is doing}
\label{sec:ch29-scope}

\subsection{Hypothesis use, statement by statement}
\label{sec:ch29-hypothesis-use}

\begin{center}
\small
\begin{tabular}{p{.28\textwidth}|p{.33\textwidth}|p{.31\textwidth}}
\toprule
\textbf{Statement}&\textbf{Uses}&\textbf{Does not use}\\
\midrule
\Cref{lem:bfj-kolodziej-sublevel},
\Cref{prop:ch29-constant-obstacle}
 &convexity of $\cE^1$; multilinearity
  \Cref{thm:e1-mixed-products}; comparison
  \eqref{eq:ch26-e1-comparison}
 &\textup{(S2)}; \textup{(S4)}; \textup{(S5)}; regularization;
  orthogonality\\
\Cref{lem:ch29-capacity-profile}
 &\Cref{prop:ch25-sublevel-capacity};
  \Cref{prop:ch25-choquet-axioms}\,\cref{itm:ch25-chq:incr};
  \Cref{lem:divisorial-positive-capacity}
 &finite energy of the potential\\
\Cref{lem:bfj-supported-domination}
 &\eqref{eq:bfj-kolodziej-sublevel};
  \Cref{lem:divisorial-positive-capacity}; continuity of $h$
 &\textup{(S5)}; \textup{(S4)}; Dini\\
\Cref{lem:ch29-iteration}
 &nothing: real analysis on $[0,\infty)$
 &all geometric hypotheses\\
\Cref{thm:ch29-kolodziej-linfty}
 &\eqref{eq:ch29-constant-obstacle};
  \eqref{eq:ch29-capacity-domination};
  \eqref{eq:ch25-normalized-bound}
 &\textup{(S5)}; orthogonality; regularization\\
\Cref{prop:ch29-skeleton-not-polar},
\Cref{prop:ch29-supported-no-polar-mass}
 &\Cref{thm:bfj-skeletal-equicontinuity}, hence \textup{(S3)}
 &capacity; energy; \textup{(S4)}\\
\Cref{prop:ch29-curve-point-capacity}
 &\Cref{prop:ch25-segment-capacities}, valid only for curves
 &anything in dimension $\geq2$\\
\Cref{thm:bfj-supported-regularity}
 &\textup{(S1)}--\textup{(S3)} through
  \Cref{thm:bfj-decreasing-regularization}; \textup{(S5)} twice;
  \textup{(S6)}; \Cref{prop:dini-usc}
 &\textup{(S4)}; orthogonality; the Ko\l odziej iteration\\
\Cref{cor:bfj-original-continuous-solution}
 &all of \textup{(S1)}--\textup{(S3)}, \textup{(S5)}, \textup{(S6)};
  orthogonality;
  \Cref{thm:ch29-uniqueness}
 &\textup{(S4)}, by
  \Cref{thm:bggjkm-remove-algebraizability}\\
\Cref{prop:ch29-model-uniqueness}
 &relative ampleness; equality clause of
  \Cref{thm:ch20-mixed-hodge}
 &capacity; energy; \textup{(S5)}\\
\Cref{thm:ch29-uniqueness}
 &\cite[Proposition~8.2]{BFJ15}; iterated Cauchy--Schwarz and ample
  Hodge equality
 &orthogonality; \textup{(S4)}; \textup{(S5)}\\
\bottomrule
\end{tabular}
\end{center}

Two rows deserve comment.  The row for
\Cref{lem:bfj-kolodziej-sublevel} says that the capacity estimates
themselves are characteristic-free: they hold as soon as the finite-energy
calculus of \Cref{chap:bfj-energy} is available, with its psh, operator and model-theoretic inputs.  Their availability
includes regularization and the relevant resolution/cofinality hypotheses;
it is not asserted to depend on residue characteristic through only
one theorem.  The row for
\Cref{thm:bfj-supported-regularity} says that \textup{(S5)} is used
\emph{twice}, once as a compactness statement in Step~1 and once as a
null-set statement in Step~2; the two uses are logically independent, and a
hypothesis supplying only one of them would not suffice.

\subsection{Later extensions}
\label{sec:ch29-later-extensions}

\emph{(E1) Positive residue characteristic.}  Every statement in
\Cref{sec:ch29-capacity} and \Cref{sec:ch29-kolodziej} survives, since
the displayed analytic deductions do not invoke multiplier ideals or
resolution anew once their psh and operator inputs are supplied.  What fails as proved
is \Cref{thm:bfj-supported-regularity}, whose first step quotes
\Cref{thm:bfj-decreasing-regularization}.  In equal characteristic $p$,
under \eqref{eq:ch29-gjkm-assumptions}, the regularization is restored by
the test-ideal theorem \Cref{thm:ch23-gjkm-positive-characteristic}, and
with it the whole chapter; this is
\Cref{thm:gjkm-positive-characteristic-cy}.  In mixed characteristic the
corresponding statement is \Cref{thm:fgk-mixed-characteristic-cy}, and it is
conditional.

\emph{(E2) Non-discrete valuations.}  The capacity estimates use only the
bounded and finite-energy Monge--Amp\`ere calculus and would transport to
any setting in which that calculus exists.  It does not exist here: over an
arbitrary complete nontrivially valued field the extension of the
Monge--Amp\`ere operator to $\cE^1$ is explicitly left open in
\cite[Remark~4.4]{PS23Hybrid}; see \Cref{prob:frontier-general-ma}.  For
curves, \Cref{thm:curves-arbitrary-field-continuity} supplies ample
metric-envelope continuity and the skeletal-supported continuous-solution
statement over every complete nontrivially valued field.  It is not by
itself the whole finite-energy package.

\emph{(E3) Singular $X$.}  Smoothness enters through the availability of
SNC models and through
\Cref{thm:bfj-skeletal-equicontinuity}, hence in
\Cref{prop:ch29-skeleton-not-polar} and in Step~1 of
\Cref{thm:bfj-supported-regularity}.  The formal estimates in \Cref{sec:ch29-capacity} transport
conditionally if the corresponding capacity and operator calculus exists.

\emph{(E4) Beyond \textup{(S5)}.} Finite-energy existence under
orthogonality is \Cref{thm:ch28-finite-energy-existence}.
The capacity-domination hypothesis gives a bounded solution once an E1
solution is available, without proving continuity.  Thuillier's curve
theory supplies another source of solutions and skeletal regularity.
The bounded nonskeletal example in
\Cref{ex:ch28-unbounded-dual-energy} demonstrates that skeletal support
is not necessary for continuity.  These statements do not constitute an
exhaustive list of regularity criteria.

\section{Dictionary}
\label{sec:ch29-dictionary}

\subsection{Objects and the statements that control them}
\label{sec:ch29-dict-objects}

\begin{center}
\small
\begin{tabular}{p{.28\textwidth}|p{.33\textwidth}|p{.31\textwidth}}
\toprule
\textbf{Analytic notion}&\textbf{Content in this chapter}
 &\textbf{Established at}\\
\midrule
sublevel set $\{\varphi<\psi\}$ of two finite-energy potentials
 &capacity controlled by a Monge--Amp\`ere mass
 &\Cref{lem:bfj-kolodziej-sublevel}\\
sublevel set $\{\varphi<-s\}$ of one potential
 &capacity profile; iteration variable
 &\Cref{prop:ch29-constant-obstacle},
  \Cref{lem:ch29-capacity-profile}\\
continuous psh minorant almost everywhere
 &minorant everywhere
 &\Cref{lem:bfj-supported-domination}\\
measure dominated by a power of capacity
 &bounded potential, explicit bound
 &\Cref{def:ch29-capacity-domination},
  \Cref{thm:ch29-kolodziej-linfty}\\
capacity floor on the support
 &capacity domination for every $\alpha$
 &\Cref{lem:ch29-capacity-floor}\\
dual complex $\Delta_{\mathcal X}$
 &disjoint from every polar set
 &\Cref{prop:ch29-skeleton-not-polar}\\
measure carried by $\Delta_{\mathcal X}$
 &no pluripolar mass; bounded pairing; weakly compact class
 &\Cref{prop:ch29-supported-no-polar-mass},
  \Cref{prop:ch29-skeletal-pairing},
  \Cref{prop:ch29-skeletal-compact}\\
decreasing model net on $\Delta_{\mathcal X}$
 &uniform convergence
 &\eqref{eq:dini-on-support}\\
finite-energy solution of a skeletal equation
 &continuous; uniform limit of model functions
 &\Cref{thm:bfj-supported-regularity},
  \Cref{cor:ch29-uniform-model-limit}\\
two solutions of the same equation
 &differ by a constant
 &\Cref{prop:ch29-model-uniqueness},
  \Cref{thm:ch29-uniqueness}\\
algebraizability over a function field
 &orthogonality in 2015; removable in 2020
 &\Cref{def:ch29-algebraizable},
  \Cref{thm:bggjkm-remove-algebraizability}\\
\bottomrule
\end{tabular}
\end{center}

\subsection{Reading the dictionary backwards}
\label{sec:ch29-dict-backward}

Five entries are one-way.  \emph{From continuity of envelopes to continuity
of solutions} there is no arrow: the former holds under the standing
hypotheses, and \Cref{ex:ch29-unbounded-not-supported} produces an
unbounded solution in its presence.  \emph{From finite energy of $\mu$ to
continuity of the solution} there is no arrow, by the same example.
\emph{From boundedness to continuity} there is no arrow proved here; see
\Cref{rem:ch29-linfty-scope}.  \emph{From support on a dual complex to
invariance under the retraction} there is no arrow, by
\cite[Theorem~B]{BJGreen24} and \Cref{warn:ch29-support-not-skeletal}.
\emph{From support on a dual complex to the capacity floor
\eqref{eq:ch29-capacity-floor}} there is a proof for the displayed
segment (\Cref{cor:ch29-curve-floor}) and for finite divisorial supports
(\Cref{rem:ch29-floor-availability}), but none is given in general
(\Cref{rem:ch29-floor-open}).

\subsection{Numbers to remember}
\label{sec:ch29-dict-numbers}

\begin{center}
\small
\begin{tabular}{p{.38\textwidth}|p{.24\textwidth}|p{.30\textwidth}}
\toprule
\textbf{Quantity}&\textbf{Value}&\textbf{Where}\\
\midrule
entry threshold of the iteration
 &$\kappa=B^{-1/\alpha}2^{-1/\alpha^2}$
 &\eqref{eq:ch29-iteration-entry}\\
total displacement of the iteration
 &$2$
 &\Cref{lem:ch29-iteration}\\
oscillation bound under $(A,\alpha)$-domination
 &$(C_\omega+nV)A^{1/\alpha}2^{n/\alpha^2}+2$
 &\eqref{eq:ch29-kolodziej-bound}\\
general capacity floor on a segment of length $\ell$
 &$\min\{m_0,1/\ell\}+\min\{m_1,1/\ell\}$
 &\eqref{eq:ch29-curve-floor}\\
that floor, and the exact infimum, for $\ell=4$, $m_0=m_1=1$
 &$1/2$ and $1$
 &\Cref{rem:ch29-floor-not-sharp}\\
$C_\omega$ for that curve
 &$4$
 &\Cref{ex:ch29-curve-solution}\\
oscillation of the uniform-measure solution
 &$1$ (bound $8$)
 &\eqref{eq:ch29-curve-solution-values},
  \eqref{eq:ch29-curve-kolodziej-numbers}\\
energy of that solution
 &$-2/3$
 &\eqref{eq:ch29-curve-energy}\\
oscillation and energy for $\mu=2\delta_{x_2}$
 &$2$ and $-2$
 &\Cref{ex:ch29-curve-dirac}\\
mass of the boundary example beyond level $m$
 &$1/(m+1)$
 &\Cref{ex:ch29-unbounded-not-supported}\\
its capacity ratio at level $m$
 &$m^{\alpha}\to\infty$
 &\eqref{eq:ch29-borderline-ratio}\\
\bottomrule
\end{tabular}
\end{center}

\section*{Source notes}

The capacity estimate, domination lemma, and continuity proof are
\cite[Lemmas~8.3--8.4 and \S8.3]{BFJ15}.  The original theorem is
\cite[Theorem~8.1]{BFJ15}.  Algebraizability is removed in
\cite{BGGJKM20}; the equal-characteristic test-ideal extension is
\cite{GJKM19}; and the conditional mixed-characteristic extension is
\cite{FGK25}.  The support and resolution hypotheses have been repeated
because they are part of the theorem, not technical decorations.

\paragraph{The original theorem, 2015.}
The sublevel estimate, domination mechanism and supported regularity
argument follow \cite[\S8]{BFJ15}.  The abstract iteration and explicit
capacity-floor deductions are worked out separately above.  The hypothesis list \eqref{eq:ch29-signature} is the one
under which Theorem~A of that paper is stated; the algebraizability
condition \textup{(S4)} is used only in
\cite[Appendix~A]{BFJ15}, and the reader may check this against the
provenance ledger \Cref{thm:orthogonality-provenance}.  The uniqueness half,
\Cref{thm:ch29-uniqueness}, is
\cite[Proposition~8.2]{BFJ15}.  Its finite-energy proof transfers
Dirichlet vanishing by iterated Cauchy--Schwarz and then uses the ample
Hodge equality, following Yuan--Zhang.  The underlying global pluripotential
theory---psh functions, compactness, envelopes, regularization---is
\cite{BFJ16}; the semipositive metrics themselves go back to
\cite{Zhang1995} and the measures to \cite{CL06}.  A short survey of the
whole circle by the same authors is \cite{BFJ16b}.

\paragraph{The proofs given here.}
\Cref{lem:bfj-kolodziej-sublevel} and
\Cref{lem:bfj-supported-domination} follow
\cite[Lemmas~8.3--8.4]{BFJ15}.  That source states skeletal support
in Lemma~8.4, but the capacity proof does not use it; the displayed
version consequently omits that unused hypothesis.  The proofs are written out in the notation of this book and, in the case
of the domination lemma, with the translation performed on \emph{both}
potentials, for the reason given in
\Cref{warn:ch29-same-translation}.  \Cref{prop:ch29-constant-obstacle},
\Cref{lem:ch29-capacity-profile}, \Cref{lem:ch29-iteration},
\Cref{lem:ch29-capacity-floor} and \Cref{thm:ch29-kolodziej-linfty} are the
non-Archimedean transcriptions of the classical Ko\l odziej argument
\cite{Kol98,Kol05}, in the form in which it is usually presented in complex
pluripotential theory; the iteration lemma with the constant
$\kappa(B,\alpha)=B^{-1/\alpha}2^{-1/\alpha^2}$ and the additive
displacement $2$ is elementary and is proved here in full so that
\Cref{ex:ch29-curve-solution} can be evaluated numerically.
No literature-priority claim is made for these explicit constants.
The bound is not needed for \Cref{thm:bfj-supported-regularity};
here it is used for the
scope discussion of \Cref{sec:ch29-support-status} and for the worked
examples.  \Cref{prop:ch29-model-uniqueness} is the equality clause of the
relative Hodge index theorem \Cref{thm:ch20-mixed-hodge} combined with
\eqref{eq:model-dirichlet-negativity}, both proved in
\Cref{chap:bfj-model-forms,chap:bfj-ma}.
\Cref{prop:ch29-skeleton-not-polar} and
\Cref{prop:ch29-supported-no-polar-mass} are immediate from
\Cref{thm:bfj-skeletal-equicontinuity} and are recorded because they are
used silently in \cite[\S8]{BFJ15}.

\paragraph{Curves, 2019.}
\Cref{thm:curves-arbitrary-field-continuity} is
based, for its ample metric-envelope half, on
\cite[Theorem~1.1]{GJKM19}, with the retraction statements
\cite[Proposition~1.2 and Propositions~3.5, 3.8]{GJKM19}; the underlying
potential theory on a Berkovich curve is Thuillier's, and the slope formula
is that of \Cref{thm:ch08-slope-formula}.  The higher-dimensional failure of
the retraction statement is the toric example of Burgos Gil and Sombra in
the appendix to \cite{GJKM19}.  The envelope half is recorded separately as
\Cref{thm:ch23-gjkm-curves} in \Cref{chap:bfj-envelopes}.

\paragraph{Equal characteristic $p$, 2019.}
\Cref{thm:gjkm-positive-characteristic-cy} is
\cite[Theorem~1.5 and Theorem~9.3]{GJKM19} under the assumptions
\eqref{eq:ch29-gjkm-assumptions}; the envelope input is
\cite[Theorem~1.4 and Theorem~8.2]{GJKM19}, recorded as
\Cref{thm:ch23-gjkm-positive-characteristic}, the orthogonality input is
\cite[Theorem~9.2]{GJKM19} together with the continuous-envelope orthogonality theorem
\cite[Theorem~6.3.2]{BGGJKM20}, the monotone regularization is
\cite[Corollary~9.1]{GJKM19}, and the uniform generation theorem for
asymptotic test ideals is
\cite[Theorem~1.3 and Theorem~4.6]{GJKM19}, generalizing
Mustaţă's theorem \cite{MustataTest} by replacing Fujita vanishing with
Keeler's generalization.  The unconditional surface case
\Cref{cor:ch29-gjkm-surface-cy} is \cite[Corollary~8.3]{GJKM19} and rests
on the three-dimensional resolution theorem of Cossart and Piltant quoted
as \cite[Theorem~6.3]{GJKM19}.

\paragraph{Removal of algebraizability, 2020.}
\Cref{thm:bggjkm-remove-algebraizability} is
\cite[Theorem~6.3.3 and Corollary~6.3.4]{BGGJKM20}; the resulting
continuous-solution statement without \textup{(S4)} is
\cite[Theorem~D]{BGGJKM20}, recorded here as
\Cref{cor:ch29-bggjkm-cy}.  The differentiability of non-Archimedean volumes
behind it is \cite[Theorem~5.4.3]{BGGJKM20}; the source also has
an appendix by Lazarsfeld.  The volume functional itself is studied in \cite{BGM21}.  As
\Cref{rem:ch29-2020-improvement} records, this work replaces one link in
the chain and leaves \textup{(S5)} untouched.

\paragraph{Mixed characteristic, 2025.}
\Cref{thm:fgk-mixed-characteristic-cy} is \cite[Theorems~7.1 and~8.1]{FGK25};
its ideal-theoretic core is \cite[Theorem~5.1]{FGK25}.  It is conditional on
two resolution properties for $K^\circ$-models, and
\Cref{rem:ch29-fgk-resolution} explains why these are not supplied by the
known three-dimensional results in mixed characteristic.  The
perturbation-friendly test ideals it uses are those of
\cite{RobertsSingh} and the surrounding literature collected in
\Cref{app:ideals}.

\paragraph{The boundary example.}
\Cref{ex:ch29-unbounded-not-supported} is assembled from
\Cref{ex:ch26-unbounded-finite-energy} of \Cref{chap:bfj-energy} and
\Cref{ex:ch25-pluripolar-computed} of \Cref{chap:bfj-capacity}; the
potential and its measure are constructed there, and the capacities are
computed there.  The observation that the resulting measure violates
\eqref{eq:ch29-capacity-domination} for every $\alpha>0$, and does so with
the exact borderline ratio \eqref{eq:ch29-borderline-ratio}, is recorded
here because it fixes the boundary of both routes to regularity at once.
The statement that a solution of a skeleton-supported equation need not be
invariant under any SNC retraction is
\cite[Theorem~B]{BJGreen24}; the accompanying fact that a divisorial
Green's function need not be piecewise $\mathbf Q$-linear is
\cite[Theorem~A]{BJGreen24}.

\paragraph{What is new here, and what is not.}
No theorem of this chapter is new.  \Cref{lem:ch29-iteration} and
\Cref{thm:ch29-kolodziej-linfty} are the standard Ko\l odziej iteration
transported to $X^{\An}$; they are written out with constants because the
scope discussion needs to know what they do and do not give.
\Cref{prop:ch29-model-uniqueness} is stated only to exhibit the mechanism
of \Cref{thm:ch29-uniqueness} in the regime where the equality clause of
the Hodge index theorem is available.  The tables of
\Cref{sec:ch29-scope,sec:ch29-dictionary} are bookkeeping, and the
computations of \Cref{sec:ch29-worked} are carried out in full because they
are the test suite against which the sign and mass conventions of
\Cref{app:normalizations} are checked in dimension one.
